\documentclass[12pt,a4paper,UTF8]{ctexart}
\usepackage[a4paper,top=22mm,bottom=22mm,left=22mm,right=22mm,headheight=15pt]{geometry}
\newcommand{\Needspace}[1]{\par\begingroup\dimen0=#1\relax\dimen2=\pagegoal\advance\dimen2 by-\pagetotal\ifdim\dimen0>\dimen2\newpage\fi\endgroup}
\newcommand{\needspace}[1]{\par\begingroup\skip0=\lastskip\vskip-\skip0\dimen0=#1\relax\advance\dimen0 by\skip0\vskip\dimen0\penalty-100\vskip-\dimen0\vskip\skip0\penalty9999\endgroup}
\usepackage{amsmath,amssymb,mathtools,bm}
\usepackage{graphicx,xcolor,booktabs,longtable,array,enumitem,fancyhdr,caption,float,microtype,setspace}
\usepackage[hidelinks,unicode]{hyperref}
\usepackage[most]{tcolorbox}
\definecolor{ContestBlue}{HTML}{17365D}
\hypersetup{pdftitle={第067套},pdfauthor={}}
\setstretch{1.24}
\setlength{\parindent}{2em}
\setlength{\parskip}{0.35em}
\setlength{\emergencystretch}{2em}
\setlist{itemsep=0.35em,topsep=0.4em}
\captionsetup{font=small,labelfont=bf,justification=centering}
\setcounter{secnumdepth}{0}
\pagestyle{fancy}
\fancyhf{}
\fancyhead[L]{\small 第067套}
\fancyfoot[C]{\thepage}
\renewcommand{\headrulewidth}{0.4pt}
\ctexset{
  section={format=\Large\bfseries\color{ContestBlue},beforeskip=1.2em,afterskip=0.7em},
  subsection={format=\large\bfseries\color{ContestBlue},beforeskip=1.0em,afterskip=0.5em},
  subsubsection={format=\normalsize\bfseries,beforeskip=0.8em,afterskip=0.35em}
}
\usepackage{calc}
\newcounter{none}
\let\originalsubsection\subsection
\renewcommand{\subsection}{\needspace{6\baselineskip}\originalsubsection}
\makeatletter
\newcommand{\tightlist}{\setlength{\itemsep}{0pt}\setlength{\parskip}{0pt}}
\newcommand{\maxwidth}{\ifdim\Gin@nat@width>\linewidth\linewidth\else\Gin@nat@width\fi}
\newcommand{\maxheight}{\ifdim\Gin@nat@height>0.72\textheight 0.72\textheight\else\Gin@nat@height\fi}
\newcommand{\pandocbounded}[1]{\begingroup\setkeys{Gin}{width=\maxwidth,height=\maxheight,keepaspectratio}#1\endgroup}
\makeatother
\begin{document}
\begin{titlepage}
  \centering
  \vspace*{58mm}
  {\Huge\bfseries\color{ContestBlue} 第067套\par}
  \vfill
  \vspace*{28mm}
\end{titlepage}
\section{第 1 题}\label{ux7b2c-1-ux9898}

光滑水平桌面上有一个哑铃：长为 \(l\) 的轻质硬杆两端分别固定质量为 \(2m\)
和 \(m\) 的小球。质量为 \(M=\lambda m\) 的小球沿 \(x\) 轴以速度 \(v\) 与
\(2m\) 球发生对心弹性碰撞。坐标原点 \(O\)
为哑铃质心，所有小球均视为刚体。若 \(M\) 球还能与哑铃发生第二次碰撞，求
\(\lambda\) 的范围。

\begin{figure}[H]
\centering
\includegraphics[width=0.648\linewidth,height=\textheight,keepaspectratio,alt={图1-1　入射小球与竖直哑铃的首次对心碰撞}]{../figures/question-01.pdf}
\caption*{图1-1　入射小球与竖直哑铃的首次对心碰撞}
\end{figure}

\subsection{一、题型定位}\label{ux4e00ux9898ux578bux5b9aux4f4d}

首次碰撞使哑铃同时平动和转动。接触点的等效质量决定冲量；第二次接触则要把小球的直线轨迹与转动杆段逐时比较。

\subsection{二、知识点}\label{ux4e8cux77e5ux8bc6ux70b9}

\begin{enumerate}
\def\labelenumi{\arabic{enumi}.}
\tightlist
\item
  一维完全弹性碰撞的速度公式。
\item
  瞬时冲量使哑铃同时获得质心平动和绕质心转动。
\item
  哑铃转动惯量需按两端到质心的实际距离计算。
\item
  在哑铃质心系中判断入射小球能否与旋转杆端再次相遇。
\item
  ``第二次碰撞''需要时间和几何可达范围同时满足。
\end{enumerate}

\subsection{三、物理图像}\label{ux4e09ux7269ux7406ux56feux50cf}

首次碰撞后，入射球相对哑铃质心向左运动，哑铃顺时针转动。哑铃转过
\(120^\circ\) 时，较长一侧的 \(m\)
端第一次到达入射球的水平相对轨迹；此后杆与该水平线的交点由最左位置向右扫过。若这一时刻入射球还没有越过杆的最左可达位置，运动球与扫动的杆必有一次相交，因而会发生第二次碰撞。

\subsection{四、解题入口}\label{ux56dbux89e3ux9898ux5165ux53e3}

先把首次碰撞暂时看成 \(M\) 与静止的 \(2m\)
球的一维弹性碰撞，求被撞端瞬时速度。再由哑铃总动量和关于质心的角动量求质心速度与角速度。最后在质心系比较入射球在哑铃转过
\(2\pi/3\) 所需时间内的位移和杆端可达到的横向距离。

\subsection{五、解析}\label{ux4e94ux89e3ux6790}

\begin{figure}[H]
\centering
\includegraphics[width=0.9\linewidth,height=\textheight,keepaspectratio,alt={图1-2　首次碰撞后哑铃质心运动、转动及二次相交几何}]{../figures/solution-01.pdf}
\caption*{图1-2　首次碰撞后哑铃质心运动、转动及二次相交几何}
\end{figure}

图1-2采用哑铃质心系；首次接触的方向等效质量为
\(2m\)；杆与小球轨迹相交还需同一时刻到达。下面先算碰撞，再检验最低可达高度。

\needspace{3\baselineskip}
\textbf{（一）首次完全弹性碰撞}

先说明为什么首次碰撞可以等效为 \(M\) 球与质量 \(2m\)
的自由质点碰撞。若接触冲量为 \(J\)，哑铃质心获得速度
\(J/(3m)\)，被撞端距质心 \(a=l/3\)，哑铃关于质心的转动惯量为
\(I_C=2ml^2/3\)。冲量还使哑铃获得角速度
\(Ja/I_C\)，所以被撞端沿冲量方向的速度增量为

\[
\Delta v_{\rm end}
=\frac{J}{3m}+\frac{Ja^2}{I_C}
=\frac{J}{3m}
+\frac{J(l/3)^2}{(2/3)ml^2}
=\frac{J}{2m}.
\]

这说明接触点沿碰撞方向所呈现的等效质量恰为
\(2m\)。与此同时，当被撞端速度为 \(v_2\) 时，哑铃的平动与转动总动能为

\[
\frac12(3m)\left(\frac23v_2\right)^2
+\frac12\left(\frac23ml^2\right)
\left(\frac{v_2}{l}\right)^2
=mv_2^2
=\frac12(2m)v_2^2.
\]

因此，把首次接触写成 \(M\) 与等效质量 \(2m\)
的一维完全弹性碰撞，并没有忽略杆的约束作用；平动和转动的效果已经包含在等效质量中。

设碰后 \(M\) 球速度为 \(v_1\)，被撞端瞬时速度为 \(v_2\)。动量守恒：

\[
Mv=Mv_1+2mv_2. \tag{1}
\]

动能守恒：

\[
\frac12Mv^2=\frac12Mv_1^2+\frac12(2m)v_2^2. \tag{2}
\]

也可使用相对速度反向关系

\[
v=v_2-v_1. \tag{3}
\]

联立并代入 \(M=\lambda m\)：

\[
\boxed{
v_1=\frac{\lambda-2}{\lambda+2}v
}, \tag{4}
\]

\[
\boxed{
v_2=\frac{2\lambda}{\lambda+2}v
}. \tag{5}
\]

\needspace{3\baselineskip}
\textbf{（二）哑铃质心位置与转动惯量}

设质心到 \(2m\) 端距离为 \(a\)，到 \(m\) 端距离为 \(b\)。有

\[
a+b=l,\qquad 2ma=mb. \tag{6}
\]

故

\[
a=\frac l3,\qquad b=\frac{2l}{3}. \tag{7}
\]

哑铃总质量为 \(3m\)。首次碰撞后总动量为 \(2mv_2\)，所以质心速度

\[
V_C=\frac{2mv_2}{3m}=\frac23v_2. \tag{8}
\]

代入式（5）：

\[
\boxed{
V_C=\frac{4\lambda}{3(\lambda+2)}v
}. \tag{9}
\]

哑铃关于质心的转动惯量

\[
\begin{aligned}
I_C
&=2m\left(\frac l3\right)^2
+m\left(\frac{2l}{3}\right)^2\\
&=\frac{2}{3}ml^2.
\end{aligned} \tag{10}
\]

碰撞冲量对质心的角动量使哑铃转动。也可直接利用被撞端速度分解：

\[
v_2=V_C+\omega\frac l3. \tag{11}
\]

所以

\[
\omega
=\frac{v_2-V_C}{l/3}
=\frac{v_2}{l}. \tag{12}
\]

代入式（5）：

\[
\boxed{
\omega=\frac{2\lambda v}{(\lambda+2)l}
}. \tag{13}
\]

\needspace{3\baselineskip}
\textbf{（三）入射球相对质心的速度}

碰后

\[
v_1-V_C
=\frac{\lambda-2}{\lambda+2}v
-\frac{4\lambda}{3(\lambda+2)}v. \tag{14}
\]

通分：

\[
\begin{aligned}
v_1-V_C
&=\frac{3\lambda-6-4\lambda}
{3(\lambda+2)}v\\
&=-\frac{\lambda+6}{3(\lambda+2)}v.
\end{aligned} \tag{15}
\]

因此 \(M\) 球在质心系中向左运动，速率为

\[
\boxed{
u=\frac{\lambda+6}{3(\lambda+2)}v
}. \tag{16}
\]

\needspace{3\baselineskip}
\textbf{（四）二次相交的最早几何时刻}

在质心系中，\(M\)
球始终沿首次碰撞点所在的水平直线运动。设杆从初始竖直方向顺时针转过
\(\phi\)。距质心 \(l/3\) 的 \(2m\) 端在该方向上的高度为

\[
y_{2m}=\frac l3\cos\phi,
\]

距质心 \(2l/3\) 的 \(m\) 端高度为

\[
y_m=-\frac{2l}{3}\cos\phi.
\]

运动球轨迹的高度为 \(y=l/3\)。当 \(0<\phi<2\pi/3\)
时，杆上各点的高度都小于 \(l/3\)，整根杆尚不能碰到该轨迹。令 \(m\)
端刚好到达轨迹，得到

\[
-\frac{2l}{3}\cos\phi=\frac l3,
\]

即

\[
\Delta\phi=\frac{2\pi}{3} \tag{17}
\]

这正是杆第一次能够重新接触运动球轨迹的时刻，而不是任意选取的``转过
\(120^\circ\)''。所需时间

\[
t_*=\frac{2\pi/3}{\omega}. \tag{18}
\]

代入式（13）：

\[
\boxed{
t_*=\frac{(\lambda+2)\pi l}{3\lambda v}
}. \tag{19}
\]

在这段时间内，入射球相对质心向左移动

\[
s_M=ut_*
=\frac{\pi(\lambda+6)}{9\lambda}l. \tag{20}
\]

此时 \(m\) 端的横坐标为

\[
x_m=-\frac{2l}{3}\sin\frac{2\pi}{3}
=-\frac{\sqrt3}{3}l.
\]

所以杆在运动球轨迹上的最左可达位置距质心

\[
s_{\max}=|x_m|=\frac{\sqrt3}{3}l. \tag{21}
\]

当 \(\phi\) 略大于 \(2\pi/3\) 时，杆与轨迹的交点从 \(x=-s_{\max}\)
向右移动，而 \(M\) 球仍向左移动。因此，只要在 \(t_*\) 时 \(M\)
球尚未越过
\(-s_{\max}\)，两者的横坐标差就会由正变负，连续性保证二者相交。故第二次碰撞的充要临界条件为

\[
s_M\le s_{\max}. \tag{22}
\]

代入式（20）、式（21）：

\[
\frac{\pi(\lambda+6)}{9\lambda}l
\le\frac{\sqrt3}{3}l. \tag{23}
\]

约去 \(l\) 并整理：

\[
\pi(\lambda+6)\le3\sqrt3\,\lambda, \tag{24}
\]

\[
6\pi\le(3\sqrt3-\pi)\lambda. \tag{25}
\]

由于 \(3\sqrt3-\pi>0\)，所以

\[
\boxed{
\lambda\ge
\frac{6\pi}{3\sqrt3-\pi}
}. \tag{26}
\]

数值上

\[
\lambda_{\min}\simeq9.18. \tag{27}
\]

等号对应 \(M\) 球恰好在 \(t_*\) 时撞到 \(m\)
端的临界情形；严格大于下限时，碰撞点一般位于随后扫过轨迹的杆段上，而不必落在端点。

\needspace{3\baselineskip}
\textbf{补充推导与适用条件}

\begin{enumerate}
\def\labelenumi{\arabic{enumi}.}
\tightlist
\item
  碰撞持续时间极短，使用冲量、动量和恢复系数。刚性哑铃并不意味着可以把杆扔掉不管；真正的理由是接触点的等效质量恰好为
  \(2m\)。这个等效质量同时正确地包含了质心平动和绕质心转动。
\item
  首次碰撞结束到第二次碰撞发生前，系统自由运动。哑铃质心做匀速直线运动，杆做匀角速度转动，\(M\)
  球也做匀速直线运动。
\item
  第二次碰撞是否发生，不再是守恒定律问题，而是相对轨迹的几何相交问题。切换到哑铃质心系后，整体平移被消去，只剩``一条水平直线''和``一根匀速扫动的杆''。
\end{enumerate}

临界不等式的方向可以从图像直接判断。\(\lambda\) 越大，重球 \(M\)
越难在第一次碰撞中改变速度，而哑铃被撞端速度趋于
\(2v\)；相对而言，杆有更多机会追上 \(M\)，所以发生二次碰撞应要求
\(\lambda\) 足够大，最终得到下限而不是上限。

还可做两个极限检查：

\[
\lambda\to\infty
\quad\Longrightarrow\quad
s_M\to\frac{\pi l}{9}
<\frac{\sqrt3l}{3},
\]

故极重的入射球确实会被再次碰到。若 \(\lambda=2\)，实验室系中 \(M\)
球碰后静止，但哑铃质心向右运动，所以质心系中 \(M\)
仍以非零速度向左运动；这说明判断二次相遇不能只盯着实验室系的 \(v_1\)。

等号情形还有明确几何意义：在杆第一次触及水平轨迹的瞬间，\(M\) 恰在 \(m\)
端的位置。大于临界值时，\(M\)
位于该点右侧，随后与向右扫过的杆段相撞；小于临界值时，\(M\)
已经跑到杆的最大几何可达范围之外，以后也不可能再相遇。

\needspace{3\baselineskip}
\textbf{杆段交点与充要条件}

图1-2用质心系，球的轨迹是 \(y=l/3\)。杆转角为 \(\phi\)
时，可将杆上点写成

\[
(x,y)=s(\sin\phi,\cos\phi),\qquad -2l/3\le s\le l/3.
\]

在 \(2\pi/3\le\phi<\pi\)，交点横坐标为 \(x_{
\rm rod}=(l/3)\tan\phi\)，从 \(-\sqrt3l/3\) 单调增加到零；球的横坐标是
\(x_M=-u\phi/\omega\)，单调减小。两者差

\[
D(\phi)=-\frac{u}{\omega}\phi-\frac l3\tan\phi,
\qquad D'(\phi)=-\frac{u}{\omega}-\frac l3\sec^2\phi<0.
\]

若 \(D(2\pi/3)\ge0\)，则到 \(\phi\to\pi^-\) 时
\(D<0\)，区间内有唯一交点；若该初值为负，则已经错过。杆在高度 \(l/3\)
的任何点都满足 \(|x|\le\sqrt3l/3\)，因为点到质心的距离不超过
\(2l/3\)。以后球只会继续向左，故晚一整周也不能补回失去的接触。这同时证明必要性与充分性。

该判据允许碰到杆段；若改题限定``只能碰两端小球''，质量范围将不再是上述整段不等式。球半径、杆厚度未给出，沿用参考解的点球、细杆几何模型。

\needspace{3\baselineskip}
\textbf{用总动能核验接触点的等效质量}

哑铃可自由平动，不能只保留绕质心转动的
\(I_C/(l/3)^2=6m\)。质心平动与转动对接触速度的贡献相加，真正的方向等效质量为

\[
m_{\rm eff}=\left(\frac1{3m}+\frac{(l/3)^2}{I_C}\right)^{-1}=2m.
\]

碰撞冲量为 \(J=2mv_2=4m\lambda v/(\lambda+2)\)，对应
\(V_C=J/(3m)\)、\(\omega=Jl/(3I_C)\)。平动和转动总动能是

\[
E_{\rm dumbbell}=\frac{J^2}{6m}+\frac{J^2l^2}{18I_C}=\frac{J^2}{4m}=mv_2^2.
\]

它与等效质点的 \(\tfrac12(2m)v_2^2\)
相同；再加碰后入射球动能，恢复碰前的
\(\tfrac12\lambda mv^2\)。若只保留转动项，便会算出错误的碰撞速度和质量比阈值。

在 \(\lambda=2\)
时入射球的实验室速度为零，但质心仍向右走，相对球便向左走；在
\(\lambda\to\infty\) 时 \(V_C\to4v/3\)、\(v_1\to v\)，相对速率趋于
\(v/3\)。两个例子都说明第二次相遇应由 \(u=V_C-v_1\) 判断，不能使用
\(v_1\)
的方向直接分类。点球与细杆模型若改成有有限尺寸的接触，还须给出半径和杆厚度。

\subsection{六、关键洞察}\label{ux516dux5173ux952eux6d1eux5bdf}

解析图采用哑铃质心系；首次接触的方向等效质量为
\(2m\)；杆与小球轨迹相交还需同一时刻到达。后续先算碰撞，再检验最低可达高度。

\subsection{七、易错点}\label{ux4e03ux6613ux9519ux70b9}

\begin{enumerate}
\def\labelenumi{\arabic{enumi}.}
\tightlist
\item
  把哑铃质心误放在杆中点。质量比为 \(2:1\)，质心到两端距离应为 \(l/3\)
  与 \(2l/3\)；一旦这里错，转动惯量、角速度和几何角度会同时出错。
\item
  直接套均匀细杆的
  \(I=3ml^2/12\)。本题杆无质量，转动惯量全部来自两个端点质点，必须逐项使用
  \(I=\sum m_ir_i^2\)。
\item
  未经说明就把碰撞对象当成孤立的 \(2m\)
  球。结论虽然恰好可用，但应通过接触点等效质量或总动能验证；否则学生会把这种做法错误推广到任意刚体端点。
\item
  把被撞端速度 \(v_2\) 当作质心速度。正确关系是
  \(v_2=V_C+\omega l/3\)，其中两项分别对应平动和转动。
\item
  用 \(v_1\) 而不是 \(v_1-V_C\)
  计算相对位移。只要比较对象是旋转哑铃，就应先把共同的质心平移消去。
\item
  认为杆转过 \(120^\circ\)
  之前已经可能撞到水平轨迹。检验杆上各点的纵坐标可知，\(2\pi/3\)
  是整根杆首次重新达到该高度的时刻。
\item
  把 \(s_M\le s_{\max}\) 理解为``在 \(t_*\)
  时必须已经碰撞''。严格地说，等号才在 \(t_*\) 碰端点；小于时是在
  \(t_*\) 之后碰到扫动的杆段。
\item
  解不等式时忽略
  \(3\sqrt3-\pi>0\)。只有确认分母为正，才能保持不等号方向。
\end{enumerate}

\subsection{八、课堂讲解建议}\label{ux516bux8bfeux5802ux8bb2ux89e3ux5efaux8bae}

讲解可按``接触点---刚体---相对轨迹''三次换视角组织。

\begin{enumerate}
\def\labelenumi{\arabic{enumi}.}
\tightlist
\item
  先不写弹性碰撞公式，只给端点一个假想冲量
  \(J\)。让学生分别算出质心速度增量和角速度增量，再把它们合成端点速度。由
  \(\Delta v_{\rm end}=J/(2m)\) 得出``端点等效质量为
  \(2m\)''，这样后面的碰撞公式有物理来源。
\item
  画出碰后速度分解图，在同一张图上标出 \(v_2\)、\(V_C\) 和
  \(\omega l/3\)。要求学生用动量和角动量各核验一次，防止把端点运动与质心运动混同。
\item
  把坐标原点移到质心，擦去实验室系的共同平移。先只画 \(M\)
  的水平轨迹，再画杆从竖直位置转到 \(2\pi/3\)
  的临界位置。让学生通过纵坐标方程
\end{enumerate}

\[
-\frac{2l}{3}\cos\phi=\frac l3
\]

自行得到 \(2\pi/3\)，而不是直接告知角度。 4. 最后把 \(s_M\) 和
\(\sqrt3l/3\)
标在同一条水平线上。可用``交点从左向右扫、球从右向左走''的连续性语言解释为何是不等式，而不是生硬记忆。

板书验收时至少保留四个中间量：

\[
V_C,\qquad \omega,\qquad u,\qquad t_*.
\]

学生若最终答案错误，可由这四项迅速定位是碰撞、刚体分解、参考系还是几何条件出了问题。

\subsection{九、变式训练}\label{ux4e5dux53d8ux5f0fux8badux7ec3}

\begin{enumerate}
\def\labelenumi{\arabic{enumi}.}
\tightlist
\item
  若入射球首先撞击质量为 \(m\)
  的长臂端，保持其他条件不变，先求该接触点的等效质量。提示：
\end{enumerate}

\[
\frac1{m_{\rm eff}}
=\frac1{3m}+\frac{(2l/3)^2}{I_C}.
\]

计算得
\(m_{\rm eff}=m\)，仍等于该端点自身质量。这是无质量杆连接两个质点、冲量垂直于杆时的性质。此时入射球轨迹的高度为
\(-2l/3\)，较短的 \(2m\)
端永远到不了该高度；只有被撞端在整周后返回。继续计算可得，点球模型中只有
\(\lambda=3\) 时入射球相对质心静止，才会在一整周后再次相遇。

\begin{enumerate}
\def\labelenumi{\arabic{enumi}.}
\setcounter{enumi}{1}
\tightlist
\item
  把完全弹性碰撞改为恢复系数 \(e\)。对
  \(1/2\le e\le1\)，仍沿用首次接触后分离的点球模型，证明
\end{enumerate}

\[
\lambda\ge\frac{12\pi e}{3\sqrt3(1+e)-2\pi(2e-1)}.
\]

提示：\(J=2m\lambda(1+e)v/(\lambda+2)\)，\(\omega=(1+e)\lambda v/[(\lambda+2)l]\)，向左相对速率为
\(u=[(2e-1)\lambda+6e]v/[3(\lambda+2)]\)。\(e=1\) 回到主答案。若
\(e<1/2\)，\(u\) 可能变号，需要重做交点分析；\(e=0\)
还须检查是否确实分离，不能把零恢复系数直接解释为黏附。

\begin{enumerate}
\def\labelenumi{\arabic{enumi}.}
\setcounter{enumi}{2}
\tightlist
\item
  将哑铃端点质量推广为 \(m_a,m_b\)，入射球撞击 \(m_a\)
  端。先证明接触点等效质量
\end{enumerate}

\[
m_{\rm eff}
=\left(\frac1{m_a+m_b}
+\frac{r_a^2}{I_C}\right)^{-1},
\]

再把本题看成 \(m_a=2m,\ m_b=m\)
的特例。该变式的重点不是算出冗长临界式，而是建立``任意刚体碰撞点都有方向等效质量''的通用认识。

\newpage

\section{第 2 题}\label{ux7b2c-2-ux9898}

长为 \(2l\) 的绝缘轻杆两端固定带电质点 \((m_1,-q_1)\) 和
\((m_2,q_2)\)，杆中点固定在无限大接地导体平板上方高度 \(h\) 处，且

\[
\frac lh=a\ll1.
\]

轻杆可绕中点自由转动。已知

\[
m_1=m(1-a),\quad q_1=q(1+a),
\]

\[
m_2=m(1+a),\quad q_2=q(1-a),
\]

以及

\[
mg=\frac{5kq^2}{8h^2},
\]

其中 \(k\) 为静电力常量。记从 \((m_1,-q_1)\) 指向 \((m_2,q_2)\)
的杆方向与向上竖直方向夹角为 \(\theta\)。

（1）求平衡角。

（2）若平衡为稳定平衡，求小振动周期。答案只含 \(m,g,h,a\) 等参数。

\begin{figure}[H]
\centering
\includegraphics[width=0.68\linewidth,height=\textheight,keepaspectratio,alt={图2-1　接地平板上方的不对称带电轻杆}]{../figures/question-02.pdf}
\caption*{图2-1　接地平板上方的不对称带电轻杆}
\end{figure}

\subsection{一、题型定位}\label{ux4e00ux9898ux578bux5b9aux4f4d-1}

导体镜像的角度势能与重力势能都是二阶小量。先保留同阶项，再由势阱曲率判断稳定性和周期；倾斜平衡角须标明近似精度。

\subsection{二、知识点}\label{ux4e8cux77e5ux8bc6ux70b9-1}

\begin{enumerate}
\def\labelenumi{\arabic{enumi}.}
\tightlist
\item
  接地无限大导体平板可用镜像电荷代替。
\item
  真实电荷与全部镜像电荷的静电能要乘 \(1/2\)。
\item
  \(a=l/h\ll1\)，只保留到决定角度依赖的 \(a^2\) 阶。
\item
  平衡由 \(dU/d\theta=0\) 决定，稳定性由 \(U''(\theta_0)>0\) 决定。
\item
  小振动周期 \(T=2\pi\sqrt{I/\kappa}\)。
\end{enumerate}

\subsection{三、物理图像}\label{ux4e09ux7269ux7406ux56feux50cf-1}

重力偏心和导体诱导的镜像电荷力矩在低阶上精细抵消，剩余角度势能为 \(a^2\)
阶。在保留 \(a^2\) 阶的近似中，倾斜的 \(\pi/3\)
是势垒顶端；两个竖直方向是曲率不同的稳定势阱。

\subsection{四、解题入口}\label{ux56dbux89e3ux9898ux5165ux53e3-1}

把两端高度写为 \(h\mp l\cos\theta\)，水平坐标写为
\(\pm l\sin\theta\)。构造重力势能和镜像电荷势能后按 \(a=l/h\) 展开。与
\(\theta\) 无关的常数可在求力矩时删除。

\subsection{五、解析}\label{ux4e94ux89e3ux6790-1}

\begin{figure}[H]
\centering
\includegraphics[width=0.88\linewidth,height=\textheight,keepaspectratio,alt={图2-2　有效角势能及三个平衡角的稳定性}]{../figures/solution-02.pdf}
\caption*{图2-2　有效角势能及三个平衡角的稳定性}
\end{figure}

图2-2左侧把实电荷与异号镜像配对，右侧按二阶角度势能画出两个势阱和中间势垒。镜像能量的
\(1/2\) 系数进入势阱曲率，从而决定小振动周期。

\needspace{3\baselineskip}
\textbf{（一）有效势能的低阶形式}

取导体板为零高度，令 \(c=\cos\theta\)。杆的方向由 \((m_1,-q_1)\) 端指向
\((m_2,q_2)\) 端，因此两端高度分别为

\[
z_1=h-lc=h(1-ac),\qquad
z_2=h+lc=h(1+ac).
\]

两端的水平间距为 \(2l\sin\theta\)。两真实电荷相距恒为
\(2l\)，所以它们之间的势能

\[
U_{12}=-\frac{kq_1q_2}{2l}
\]

与 \(\theta\) 无关，可以并入常数
\(U_0\)，但不能把导体诱导电荷产生的角度依赖也一并删去。

先算重力势能：

\[
\begin{aligned}
U_g
&=m_1gz_1+m_2gz_2\\
&=mgh\left[(1-a)(1-ac)+(1+a)(1+ac)\right]\\
&=2mgh+2mgh\,a^2c.
\end{aligned}
\]

这里一阶项消失，是质量不对称和几何小量共同作用的结果。

真实电荷 \(-q_1\)、\(q_2\) 在导体板下方的镜像分别为
\(+q_1\)、\(-q_2\)。镜像法求力时可以直接使用镜像电荷；求系统静电能时，由于诱导电荷是随真实电荷逐步建立的，必须取

\[
U_{\rm im}=\frac12\sum_i q_i\varphi_{\rm image}(\boldsymbol r_i).
\]

因此两个电荷与各自镜像的能量为

\[
U_{\rm self}
=-\frac{kq_1^2}{4z_1}-\frac{kq_2^2}{4z_2}.
\]

两个``真实电荷---对方镜像''的距离相同，均为

\[
\begin{aligned}
d
&=\sqrt{(2l\sin\theta)^2+(z_1+z_2)^2}\\
&=2h\sqrt{1+a^2\sin^2\theta}.
\end{aligned}
\]

交叉项在求和后为

\[
U_{\rm cross}=\frac{kq_1q_2}{d}.
\]

于是

\[
\begin{aligned}
U_{\rm self}
&=-\frac{kq^2}{4h}
\left[
\frac{(1+a)^2}{1-ac}
+\frac{(1-a)^2}{1+ac}
\right],\\
U_{\rm cross}
&=\frac{kq^2}{2h}
\frac{1-a^2}{\sqrt{1+a^2(1-c^2)}}.
\end{aligned}
\]

下面把每个因子保留到 \(a^2\)。使用

\[
\frac1{1\mp ac}=1\pm ac+a^2c^2+O(a^4),
\]

以及

\[
\frac1{\sqrt{1+a^2(1-c^2)}}
=1-\frac12a^2(1-c^2)+O(a^4),
\]

可得

\[
\begin{aligned}
\frac{(1+a)^2}{1-ac}
+\frac{(1-a)^2}{1+ac}
&=2+2a^2(1+2c+c^2)+O(a^4),\\
\frac{1-a^2}{\sqrt{1+a^2(1-c^2)}}
&=1-\frac32a^2+\frac12a^2c^2+O(a^4).
\end{aligned}
\]

代回并只保留含 \(c\) 的项：

\[
\begin{aligned}
U_{\rm self}+U_{\rm cross}
&=U_{\rm e0}
-\frac{kq^2}{h}a^2c
-\frac{kq^2}{4h}a^2c^2
+O(a^4).
\end{aligned}
\]

再与重力势能相加，角度相关部分为

\[
U_{\rm ang}
=a^2\left(2mgh-\frac{kq^2}{h}\right)c
-\frac{kq^2}{4h}a^2c^2.
\]

利用题给参数关系

\[
mg=\frac{5kq^2}{8h^2}, \tag{1}
\]

即

\[
\frac{kq^2}{h}=\frac85mgh,
\]

逐项化简：

\[
2mgh-\frac{kq^2}{h}
=\left(2-\frac85\right)mgh
=\frac25mgh,
\]

\[
-\frac{kq^2}{4h}
=-\frac14\cdot\frac85mgh
=-\frac25mgh.
\]

故所有与角度有关的项最终化为

\[
\boxed{
U_{\rm eff}(\theta)
=U_0+\frac25mgh\,a^2
\left(\cos\theta-\cos^2\theta\right)
}. \tag{2}
\]

\(U_0\) 包含两真实电荷间的常量势能及镜像展开中与 \(\theta\)
无关的项，不影响平衡与小振动。

量纲检查：\(mgh\) 是能量，\(a\) 及三角函数无量纲，故式（2）量纲正确。把
\(a\to0\)
时，角度势能消失，也符合两端趋于同一点后方向不再有物理意义的极限。

\needspace{3\baselineskip}
\textbf{（二）求平衡角}

对式（2）求导：

\[
\frac{dU_{\rm eff}}{d\theta}
=\frac25mgh\,a^2
\left(-\sin\theta+2\sin\theta\cos\theta\right). \tag{3}
\]

提取公因子：

\[
\frac{dU_{\rm eff}}{d\theta}
=\frac25mgh\,a^2
\sin\theta(2\cos\theta-1). \tag{4}
\]

平衡条件为

\[
\sin\theta(2\cos\theta-1)=0. \tag{5}
\]

在 \(0\le\theta\le\pi\) 中，

\[
\boxed{
\theta=0,\quad \frac{\pi}{3},\quad \pi
}. \tag{6}
\]

\needspace{3\baselineskip}
\textbf{（三）判断稳定性}

系统所受回复力矩

\[
\mathcal M(\theta)
=-\frac{dU_{\rm eff}}{d\theta}
=-\frac25mgh\,a^2
\sin\theta(2\cos\theta-1). \tag{7}
\]

在 \(\theta_0=0\) 附近令 \(\theta=\delta\)：

\[
\sin\theta\simeq\delta,\qquad
2\cos\theta-1\simeq1. \tag{8}
\]

所以

\[
\mathcal M\simeq-\frac25mgh\,a^2\delta, \tag{9}
\]

是回复力矩，故 \(\theta=0\) 稳定。

在 \(\theta_0=\pi/3\) 附近，力矩线性项与位移同向，是不稳定平衡。

在 \(\theta_0=\pi\) 附近令 \(\theta=\pi+\delta\)：

\[
\sin\theta\simeq-\delta,\qquad
2\cos\theta-1\simeq-3. \tag{10}
\]

故

\[
\mathcal M\simeq-\frac65mgh\,a^2\delta, \tag{11}
\]

也是回复力矩。因此

\[
\boxed{\theta=0,\pi\text{ 稳定；}\ \theta=\pi/3\text{ 不稳定}}. \tag{12}
\]

\needspace{3\baselineskip}
\textbf{（四）转动惯量}

两质点到转轴距离均为 \(l\)：

\[
I=(m_1+m_2)l^2. \tag{13}
\]

而

\[
m_1+m_2=m(1-a)+m(1+a)=2m, \tag{14}
\]

所以

\[
I=2ml^2=2ma^2h^2. \tag{15}
\]

\needspace{3\baselineskip}
\textbf{（五）两个稳定位置的小振动周期}

在 \(\theta=0\) 附近，等效回复力矩系数

\[
\kappa_0=\frac25mgh\,a^2. \tag{16}
\]

因此

\[
\begin{aligned}
T_0
&=2\pi\sqrt{\frac{I}{\kappa_0}}\\
&=2\pi
\sqrt{
\frac{2ma^2h^2}
{(2/5)mgh\,a^2}
}\\
&=\boxed{2\pi\sqrt{\frac{5h}{g}}}.
\end{aligned} \tag{17}
\]

在 \(\theta=\pi\) 附近，

\[
\kappa_\pi=\frac65mgh\,a^2. \tag{18}
\]

故

\[
\boxed{
T_\pi=2\pi\sqrt{\frac{5h}{3g}}
}. \tag{19}
\]

参数 \(m,a\) 均消去，是因为转动惯量和势阱曲率都正比于 \(ma^2\)。

\needspace{3\baselineskip}
\textbf{补充推导与适用条件}

\[
mg=\frac{5kq^2}{8h^2}
\]

后，系数被精确组合成同一尺度 \((2/5)mgh\,a^2\)，最终只剩

\[
U_{\rm eff}-U_0
=A(c-c^2),\qquad
A=\frac25mgh\,a^2,\quad c=\cos\theta.
\]

一旦得到这个形式，复杂的静电问题就转化为一元函数图像问题。平衡条件既可以对
\(\theta\) 求导，也可以先看抛物线 \(c-c^2\)；但因为
\(c=\cos\theta\)，端点 \(c=\pm1\) 对应的 \(\sin\theta=0\)
也必须保留，不能只解 \(dU/dc=0\)。

镜像能量中的 \(1/2\)
是本题最本质的静电细节。镜像电荷只是计算导体响应的数学替身，不是另一组可以独立搬运到无穷远的真实电荷。若直接把所有真实---镜像库仑势能相加，能量会多算一倍，平衡角和周期都会改变。

两个稳定点的曲率不同：

\[
U''(0)=\frac25mgh\,a^2,\qquad
U''(\pi)=\frac65mgh\,a^2.
\]

因此 \(\theta=\pi\) 附近的势阱更``硬''，周期比 \(\theta=0\) 附近短
\(\sqrt3\) 倍。势能值也不同：

\[
U(0)=U_0,\qquad
U(\pi)=U_0-\frac45mgh\,a^2,
\]

所以有微弱耗散时，系统最终更倾向落入 \(\theta=\pi\) 的较深势阱。

小振动周期中 \(a^2\)
消失不是偶然。势阱曲率随杆长平方减小，而转动惯量同样随杆长平方减小；两者之比只留下重力时间尺度
\(\sqrt{h/g}\)。这也是答案量纲检查的最快方法。

\needspace{3\baselineskip}
\textbf{二阶近似与有限杆长的区别}

前面的周期为 \(a\ll1\) 的首项，倾斜角是
\(\theta=\pi/3+O(a^2)\)。把角度势能除以
\(mgh\)，不含角度的真实电荷互作用仍可删去，精确表达式为

\[
u(c)=2a^2c-\frac25\left[\frac{(1+a)^2}{1-ac}+\frac{(1-a)^2}{1+ac}\right]
+\frac45\frac{1-a^2}{\sqrt{1+a^2(1-c^2)}},\quad c=\cos\theta.
\]

交换 \(a\mapsto-a\) 不改变这条式子，所以下一个误差是 \(a^4\)，不是
\(a^3\)。其展开为

\[
u(c)=-2a^2+\frac25a^2(c-c^2)+a^4 f(c)+O(a^6),
\]

\[
f(c)=-\frac45c^4-\frac85c^3-\frac65c^2+\frac3{10}(1-c^2)^2+\frac25.
\]

\(f'(1/2)=-13/4\)，所以倾斜驻点有

\[
\cos\theta_* =\frac12-\frac{65}{16}a^2+O(a^4),\qquad
\theta_* =\frac\pi3+\frac{65}{8\sqrt3}a^2+O(a^4).
\]

图2-2画的是二阶近似曲线，不能拿图上的 \(\pi/3\) 作为任意 \(a\)
的精确角。两个竖直驻点则由轴对称性保持在 \(0,\pi\)。直接微分精确势能可得

\[
\frac{U''(0)}{mgh}=\frac{2a^2(2a^6-11a^4+24a^2+1)}{5(1-a^2)^2},
\quad
\frac{U''(\pi)}{mgh}=\frac{2a^2(2a^2+3)}5.
\]

除以 \(I=2ma^2h^2\) 即得有限 \(a\) 的小振动频率；再令
\(a\to0\)，恢复前述两个周期。若允许三维转动，在竖直位置两个互相垂直的微小倾斜具有相同频率；倾斜驻点的方位角方向中性，不能把方位自由度也称为回复振动。

\needspace{3\baselineskip}
\textbf{两种周期为何相差 \(\sqrt3\)}

令 \(K=2mgh\,a^2/5\)，二阶角势能的变量部分为
\(K(\cos\theta-\cos^2\theta)\)。对两个竖直状态分别展开：

\[
U(\vartheta)=U(0)+\frac K2\vartheta^2+O(\vartheta^4),
\]

\[
U(\pi+\vartheta)=U(\pi)+\frac{3K}{2}\vartheta^2+O(\vartheta^4).
\]

杆系绕中点的转动惯量为 \(I=2mh^2a^2\)，所以
\(\omega_0^2=K/I=g/(5h)\)，\(\omega_\pi^2=3K/I=3g/(5h)\)。相同惯量配合三倍势阱曲率，便给出
\(T_0/T_\pi=\sqrt3\)。

倾斜平衡点不是第三个势阱。二阶式的导数为
\(K\sin\theta(2\cos\theta-1)\)，在 \(\theta=\pi/3\) 处二阶导数为
\(-3K/2\)，故小偏离会远离该点。若允许杆在水平面内改变方位角，势能只依赖与竖直的夹角；方位方向并没有同样的恢复曲率，不能为它再写一个有限小振动周期。

\subsection{六、关键洞察}\label{ux516dux5173ux952eux6d1eux5bdf-1}

解析图左侧把实电荷与异号镜像配对，右侧按二阶角度势能画出两个势阱和中间势垒。镜像能量的
\(1/2\) 系数进入势阱曲率，从而决定小振动周期。

\subsection{七、易错点}\label{ux4e03ux6613ux9519ux70b9-1}

\begin{enumerate}
\def\labelenumi{\arabic{enumi}.}
\tightlist
\item
  把杆两端标签放反。\(\theta\) 是从由 \(m_1\) 指向 \(m_2\)
  的杆方向量到竖直方向；因此
  \(z_1=h-l\cos\theta,\ z_2=h+l\cos\theta\)。若高度写反，重力项的一次系数会变号。
\item
  镜像电荷符号写错。接地平面镜像必须与对应真实电荷异号，所以 \(-q_1\)
  的镜像是 \(+q_1\)，\(+q_2\) 的镜像是 \(-q_2\)。
\item
  计算静电能时漏掉 \(1/2\)。自镜像项应为 \(-kq_i^2/(4z_i)\)，不是
  \(-kq_i^2/(2z_i)\)。
\item
  交叉镜像项重复或漏计。两个真实电荷分别与对方镜像作用，在总和前有
  \(1/2\)，最终恰好留下一个 \(kq_1q_2/d\)。
\item
  只把分母展开而不展开 \(q_1^2,q_2^2\)。\((1\pm a)^2\) 中的 \(a^2\)
  与几何展开同阶，必须同时保留。
\item
  看到一阶项抵消就停止，误判势能与角度无关。本题的第一个非零角度项正是
  \(a^2\) 阶。
\item
  只解 \(2\cos\theta-1=0\)，漏掉 \(\sin\theta=0\) 给出的
  \(\theta=0,\pi\)。
\item
  只凭``势能大或小''口头判断稳定性。应计算 \(U''(\theta_0)\)，或把
  \(\theta=\theta_0+\delta\) 代入力矩并检查是否为 \(-\kappa\delta\)。
\item
  转动惯量误用杆全长 \(2l\)。两个质点到轴的距离都是 \(l\)，所以
  \(I=(m_1+m_2)l^2\)。
\end{enumerate}

\subsection{八、课堂讲解建议}\label{ux516bux8bfeux5802ux8bb2ux89e3ux5efaux8bae-1}

建议把板书分成``几何---能量---展开---曲率''四列。

\begin{enumerate}
\def\labelenumi{\arabic{enumi}.}
\tightlist
\item
  几何列只写
  \(z_1,z_2,d\)。要求学生说明每个距离对应哪一对电荷，尤其辨认交叉镜像距离中的竖直分量是
  \(z_1+z_2=2h\)。
\item
  能量列分成 \(U_g,U_{12},U_{\rm self},U_{\rm cross}\)。先用颜色标出
  \(U_{12}\) 与角度无关，再解释为何它虽然很大，却不产生力矩。
\item
  展开列先单独列出
\end{enumerate}

\[
(1\pm x)^{-1},\qquad (1+x)^{-1/2}
\]

的二阶公式。每完成一个展开就按``常数、\(c\)、\(c^2\)''归类，避免在一行长式子中同时展开、合并和代入条件。
4. 曲率列从统一的 \(U_{\rm eff}\) 出发。让学生分别在三个平衡角附近写
\(\theta=\theta_0+\delta\)，从力矩的线性项辨认稳定性，再用
\(I\ddot\delta=-\kappa\delta\) 得周期。

课堂验收可设置三个快速问题：为什么真实电荷间势能不影响平衡？为什么镜像能要乘
\(1/2\)？为什么周期不含
\(a\)？若学生能分别用``距离恒定''\,``诱导电荷逐步建立''\,``惯量与曲率同阶''回答，说明已经掌握而不是只会抄公式。

图2-2应在推导完式（2）后使用：沿曲线指认 \(\theta=0,\pi/3,\pi\)
三点，比较斜率为零与曲率正负。不要在未建立势能前先给稳定结论。

\subsection{九、变式训练}\label{ux4e5dux53d8ux5f0fux8badux7ec3-1}

\begin{enumerate}
\def\labelenumi{\arabic{enumi}.}
\tightlist
\item
  在有效势能式（2）的基础上画出 \(0\le\theta\le\pi\)
  的曲线，计算三个平衡点的势能和曲率。若从 \(\theta=0\)
  处给系统一个角速度 \(\Omega\)，求越过 \(\theta=\pi/3\) 势垒所需的最小
  \(\Omega\)。提示使用
\end{enumerate}

\[
\frac12I\Omega_{\min}^2
=U\!\left(\frac{\pi}{3}\right)-U(0).
\]

\begin{enumerate}
\def\labelenumi{\arabic{enumi}.}
\setcounter{enumi}{1}
\tightlist
\item
  将题给参数关系改成
\end{enumerate}

\[
mg=\beta\frac{kq^2}{h^2},
\]

保留其他不变。把有效势能写成

\[
U_{\rm ang}=a^2(A\cos\theta-B\cos^2\theta),
\]

求 \(A,B\) 关于 \(\beta\) 的表达式，并讨论倾斜平衡存在的条件
\(|A/(2B)|\le1\)。

\begin{enumerate}
\def\labelenumi{\arabic{enumi}.}
\setcounter{enumi}{2}
\tightlist
\item
  若杆本身有均匀质量 \(m_r\)，静电势能不变，但转动惯量增加
\end{enumerate}

\[
I_r=\frac1{12}m_r(2l)^2.
\]

重新求两个稳定位置的小振动周期。该变式用于强调：平衡角只由势能决定，周期还取决于惯量。

\newpage

\section{第 3 题}\label{ux7b2c-3-ux9898}

两块水平固定隔板间距为 \(H_0\)。两质量均为 \(m\)
的振子由三根竖直弹簧串接在隔板之间，三根弹簧的劲度系数依次为
\(k,K,k\)，原长依次为 \(l_0,L_0,l_0\)，且

\[
2l_0+L_0=H_0.
\]

取上隔板处为原点、竖直向下为正。

（1）求平衡时三根弹簧从上到下的形变量。

（2）系统已平衡，给上方振子一个向下初速度
\(v_0\)，下方振子初速度为零。求两个振子的绝对位置随时间的函数。

\begin{figure}[H]
\centering
\includegraphics[width=0.282\linewidth,height=\textheight,keepaspectratio,alt={图3-1　两振子与三弹簧串联系统及坐标约定}]{../figures/question-03.pdf}
\caption*{图3-1　两振子与三弹簧串联系统及坐标约定}
\end{figure}

\subsection{一、题型定位}\label{ux4e00ux9898ux578bux5b9aux4f4d-2}

重力先改变平衡位置，随后的小振动由弹力增量决定。相加与相减可分离同相、反相模态，最后还要把相对位移加回绝对位置。

\subsection{二、知识点}\label{ux4e8cux77e5ux8bc6ux70b9-2}

\begin{enumerate}
\def\labelenumi{\arabic{enumi}.}
\tightlist
\item
  固定端距等于三弹簧原长之和，所以三个形变量代数和为零。
\item
  关于平衡位置写动力学方程时，重力项自动消失。
\item
  对称系统的同相、反相组合可直接解耦。
\item
  初速度激励通常同时激发两个正规模态。
\item
  最终答案要加回平衡位置，不能只给相对位移。
\end{enumerate}

\subsection{三、物理图像}\label{ux4e09ux7269ux7406ux56feux50cf-2}

同相模态中两振子一起上下移动，中间弹簧长度不变，频率只由外侧弹簧 \(k\)
决定。反相模态中两振子相向或相背运动，中间弹簧形变加倍，频率升高到
\(\sqrt{(k+2K)/m}\)。

\subsection{四、解题入口}\label{ux56dbux89e3ux9898ux5165ux53e3-2}

先用形变量总和与两个振子的静力平衡求平衡位置。动力学部分令 \(x_a,x_b\)
为相对各自平衡点的向下位移，写出耦合方程，再定义

\[
q_1=x_a+x_b,\qquad q_2=x_a-x_b.
\]

\subsection{五、解析}\label{ux4e94ux89e3ux6790-2}

\begin{figure}[H]
\centering
\includegraphics[width=0.86\linewidth,height=\textheight,keepaspectratio,alt={图3-2　同相模态与反相模态的位移方向}]{../figures/solution-03.pdf}
\caption*{图3-2　同相模态与反相模态的位移方向}
\end{figure}

图3-2把两个物块相对平衡位置的位移分为同相与反相。弹簧长度先用于求静力平衡，随后只对位移增量列运动方程。

\needspace{3\baselineskip}
\textbf{（一）静力平衡}

记从上到下三根弹簧的形变量为
\(\Delta l_1,\Delta l_2,\Delta l_3\)，伸长取正。固定端距不变且等于原长总和：

\[
\Delta l_1+\Delta l_2+\Delta l_3=0. \tag{1}
\]

对上方振子列竖直方向平衡：

\[
k\Delta l_1=mg+K\Delta l_2. \tag{2}
\]

对下方振子：

\[
K\Delta l_2=mg+k\Delta l_3. \tag{3}
\]

由式（2）

\[
\Delta l_1=\frac{mg}{k}+\frac{K}{k}\Delta l_2, \tag{4}
\]

由式（3）

\[
\Delta l_3=\frac{K}{k}\Delta l_2-\frac{mg}{k}. \tag{5}
\]

代入式（1）：

\[
\left(1+\frac{2K}{k}\right)\Delta l_2=0. \tag{6}
\]

故

\[
\boxed{
\Delta l_1=\frac{mg}{k},\quad
\Delta l_2=0,\quad
\Delta l_3=-\frac{mg}{k}
}. \tag{7}
\]

下弹簧处于压缩状态。

上振子的平衡坐标

\[
X_{a0}=l_0+\frac{mg}{k}. \tag{8}
\]

下振子的平衡坐标

\[
X_{b0}=H_0-\left(l_0-\frac{mg}{k}\right)
=H_0-l_0+\frac{mg}{k}. \tag{9}
\]

\needspace{3\baselineskip}
\textbf{（二）关于平衡位置的动力学方程}

设

\[
X_a=X_{a0}+x_a,\qquad
X_b=X_{b0}+x_b. \tag{10}
\]

重力与静态弹力已经在平衡位置抵消。上振子受增量弹力：

\[
m\ddot x_a
=-kx_a+K(x_b-x_a). \tag{11}
\]

即

\[
m\ddot x_a+(k+K)x_a-Kx_b=0. \tag{12}
\]

下振子：

\[
m\ddot x_b
=-kx_b-K(x_b-x_a), \tag{13}
\]

即

\[
m\ddot x_b+(k+K)x_b-Kx_a=0. \tag{14}
\]

\needspace{3\baselineskip}
\textbf{（三）正规坐标解耦}

定义

\[
q_1=x_a+x_b,\qquad q_2=x_a-x_b. \tag{15}
\]

式（12）与式（14）相加：

\[
\ddot q_1+\frac{k}{m}q_1=0. \tag{16}
\]

所以

\[
\boxed{\omega_1=\sqrt{\frac{k}{m}}}. \tag{17}
\]

两式相减：

\[
\ddot q_2+\frac{k+2K}{m}q_2=0. \tag{18}
\]

所以

\[
\boxed{\omega_2=\sqrt{\frac{k+2K}{m}}}. \tag{19}
\]

\needspace{3\baselineskip}
\textbf{（四）代入初始条件}

\(t=0\) 时系统在平衡位置：

\[
x_a(0)=x_b(0)=0. \tag{20}
\]

初速度为

\[
\dot x_a(0)=v_0,\qquad
\dot x_b(0)=0. \tag{21}
\]

因此

\[
q_1(0)=q_2(0)=0, \tag{22}
\]

\[
\dot q_1(0)=\dot q_2(0)=v_0. \tag{23}
\]

所以

\[
q_1(t)=\frac{v_0}{\omega_1}\sin\omega_1t, \tag{24}
\]

\[
q_2(t)=\frac{v_0}{\omega_2}\sin\omega_2t. \tag{25}
\]

由定义反解：

\[
x_a=\frac{q_1+q_2}{2},\qquad
x_b=\frac{q_1-q_2}{2}. \tag{26}
\]

最终

\[
\boxed{
\begin{aligned}
X_a(t)
&=l_0+\frac{mg}{k}
+\frac{v_0}{2\omega_1}\sin\omega_1t
+\frac{v_0}{2\omega_2}\sin\omega_2t,\\
X_b(t)
&=H_0-l_0+\frac{mg}{k}
+\frac{v_0}{2\omega_1}\sin\omega_1t
-\frac{v_0}{2\omega_2}\sin\omega_2t,
\end{aligned}
}. \tag{27}
\]

其中

\[
\omega_1=\sqrt{\frac{k}{m}},
\qquad
\omega_2=\sqrt{\frac{k+2K}{m}}. \tag{28}
\]

\needspace{3\baselineskip}
\textbf{补充推导与适用条件}

动力学部分的本质是对称性。关于平衡点的方程可写成

\[
m\ddot{\boldsymbol x}
+\begin{pmatrix}
k+K&-K\\
-K&k+K
\end{pmatrix}
\boldsymbol x=0.
\]

刚度矩阵在交换上下两个振子时不变，所以它的本征向量必可取为交换操作的两个本征方向：

\[
\boldsymbol e_1=(1,1),\qquad
\boldsymbol e_2=(1,-1).
\]

这就是同相与反相模态。正规坐标不是人为猜测出来的代数花招，而是由系统的镜像对称性决定。

同相时 \(x_b-x_a=0\)，中间弹簧没有增量形变，所以 \(K\) 不进入
\(\omega_1\)。反相时
\(x_b-x_a=-2x_a\)，中间弹簧对每个振子的回复效应加倍，所以刚度从 \(k\)
增加到 \(k+2K\)。这一图像可以在写矩阵之前预测两个频率。

初态只给上振子速度，相当于速度向量

\[
(v_0,0)
=\frac{v_0}{2}(1,1)
+\frac{v_0}{2}(1,-1),
\]

因此两个模态被等量激发。这里``等量''指模态速度投影相等，不是位移振幅相等；振幅还要除以各自频率，所以低频同相模态的位移振幅更大。

最终答案可做三重核验。代入
\(t=0\)，两位移应回到各自平衡坐标；对时间求导，必须得到
\(\dot X_a(0)=v_0,\dot X_b(0)=0\)；再求二阶导并代回式（12）、式（14），两种频率项应分别独立满足方程。

极限 \(K\to0\) 时
\(\omega_2\to\omega_1\)，两个振子退化为相同频率。此时式（27）化为上振子单独振动、下振子保持静止，正好符合中间弹簧消失后的物理图像。

\needspace{3\baselineskip}
\textbf{能量与实际可运动区间}

相对平衡位置的机械能为

\[
E=\frac m2(\dot x_a^2+\dot x_b^2)+\frac k2(x_a^2+x_b^2)+\frac K2(x_b-x_a)^2.
\]

代入 \(x_a=(q_1+q_2)/2\)、\(x_b=(q_1-q_2)/2\)，得到

\[
E=\frac m4(\dot q_1^2+\omega_1^2q_1^2)
+\frac m4(\dot q_2^2+\omega_2^2q_2^2).
\]

初始两项各为 \(mv_0^2/4\)，总能量为
\(mv_0^2/2\)；两模态的能量相等，位移振幅却分别是 \(v_0/\omega_1\) 和
\(v_0/\omega_2\)。图3-2的两组箭头分别表示同相、反相的瞬时位移组合，不能当成始终固定方向的速度或弹力。

静态下弹簧长度为 \(l_0-mg/k\)，须为正。运动时三段长度分别为

\[
\begin{aligned}
\ell_1&=l_0+mg/k+x_a,\\
\ell_2&=L_0+x_b-x_a,\\
\ell_3&=l_0-mg/k-x_b.
\end{aligned}
\]

题中理想弹簧解需在这三段仍保持正长、没有碰到隔板的范围内使用。\(|x_a|,|x_b|\le (v_0/2)(1/\omega_1+1/\omega_2)\)
给出简便的充分幅度约束。若真实弹簧只承受拉力，压缩支路的模型会改变，不能直接套同一解。

\needspace{3\baselineskip}
\textbf{单边初速为何激发两种模态}

初始 \(x_a=x_b=0\)、\(\dot x_a=v_0\)、\(\dot x_b=0\)，故
\(\dot q_1=\dot q_2=v_0\)。两个广义坐标分别为

\[
q_1=\frac{v_0}{\omega_1}\sin\omega_1t,\qquad
q_2=\frac{v_0}{\omega_2}\sin\omega_2t.
\]

由它们复原后，两物块速度为

\[
\dot x_a=\frac{v_0}{2}(\cos\omega_1t+\cos\omega_2t),
\]

\[
\dot x_b=\frac{v_0}{2}(\cos\omega_1t-\cos\omega_2t).
\]

在 \(t=0\) 处应分别还原 \(v_0,0\)；将位移代回两条运动方程，又应分别还原
\(-kx_a-K(x_a-x_b)\) 和
\(-kx_b-K(x_b-x_a)\)。这比只检查两个频率更能发现符号或模态系数错误。

当 \(K\to0\)，两个频率重合，\(x_b\)
恒为零；这符合中间弹簧不传递作用的极限。当
\(K\gg k\)，反相模态频率变大、位移振幅变小，但它的初始能量仍为
\(mv_0^2/4\)，不能仅凭振幅小就删除这一模态。

\subsection{六、关键洞察}\label{ux516dux5173ux952eux6d1eux5bdf-2}

解析图把两个物块相对平衡位置的位移分为同相与反相。弹簧长度先用于求静力平衡，随后只对位移增量列运动方程。

\subsection{七、易错点}\label{ux4e03ux6613ux9519ux70b9-2}

\begin{enumerate}
\def\labelenumi{\arabic{enumi}.}
\tightlist
\item
  认为三根弹簧在重力作用下都伸长。下端固定，向下的物块会压缩最下方弹簧；``伸长取正''的符号下应有
  \(\Delta l_3<0\)。
\item
  对下振子列平衡时把下弹簧弹力方向写反。用有符号形变量统一处理最稳妥：式（3）解出
  \(\Delta l_3=-mg/k\)，负号自然表示压缩。
\item
  漏掉总长度约束
  \(\Delta l_1+\Delta l_2+\Delta l_3=0\)，导致静力方程不封闭。
\item
  关于平衡点写增量方程后又重复加入
  \(mg\)。静态弹力已经抵消重力，\(x_a,x_b\) 方程中只保留弹力增量。
\item
  中间弹簧增量写成
  \(x_a+x_b\)。两物块都向下移动相同距离时，中间弹簧长度不变，所以正确形式必须在
  \(x_a=x_b\) 时为零，即 \(x_b-x_a\)。
\item
  相加、相减方程时漏掉质量因子，或把反相频率写成
  \(\sqrt{(k+K)/m}\)。中间弹簧在反相模态中的长度变化是单个振幅的两倍。
\item
  把 \(q_1,q_2\) 直接当成某个物块位移。反解时有
\end{enumerate}

\[
x_a=\frac{q_1+q_2}{2},\qquad
x_b=\frac{q_1-q_2}{2},
\]

因子 \(1/2\) 不能漏。 8. 只写相对平衡位置的
\(x_a,x_b\)，没有回答题目要求的绝对位置 \(X_a,X_b\)。 9.
初速度投影时误写 \(\dot q_2(0)=-v_0\)。由 \(q_2=x_a-x_b\)，本题应为
\(+v_0\)。

\subsection{八、课堂讲解建议}\label{ux516bux8bfeux5802ux8bb2ux89e3ux5efaux8bae-2}

先讲静态，再讲动态，且明确两套坐标不能混用。

\begin{enumerate}
\def\labelenumi{\arabic{enumi}.}
\tightlist
\item
  在题图旁画三个带正负号的形变量箭头。先让学生只用``固定端距离不变''写出代数和为零，再分别隔离两个物块画受力图。求出下弹簧压缩后，追问为什么这不与弹簧只能``拉''相矛盾：理想螺旋弹簧既能拉伸也能压缩。
\item
  求得平衡坐标后，另起一栏定义
  \(x_a,x_b\)。用彩色括号标出静态力与增量力，说明重力为何从增量方程中消失。
\item
  不急于做代数相加减，先展示图3-2。让学生观察同相运动时中间弹簧不变、反相运动时中间弹簧变化两倍，并据此预测
  \(\omega_1,\omega_2\) 的结构。
\item
  再把预测与矩阵或相加减法相互验证。对刚接触竞赛的学生，相加减法是主线；矩阵只作为解释对称性的补充，不要求先掌握行列式。
\item
  初始条件部分用速度向量分解：
\end{enumerate}

\[
(v_0,0)
=\frac12(v_0,v_0)+\frac12(v_0,-v_0).
\]

这一行可以直接解释为何两种模态同时出现。 6.
结尾必须现场核验四个初值。把最终两式分别代入 \(t=0\)
并求导，学生会直观看到两个模态在下振子初速度中相消、在上振子初速度中相加。

若课堂时间允许，可播放或手动画出 \(K\gg k\)
时的运动：反相模态很快、同相模态很慢，叠加后会出现快振动叠在慢摆动上的波形。这比只报两个根更能建立双频直觉。

\subsection{九、变式训练}\label{ux4e5dux53d8ux5f0fux8badux7ec3-2}

\begin{enumerate}
\def\labelenumi{\arabic{enumi}.}
\tightlist
\item
  若两振子初始速度均为 \(v_0\)，且初位移为零，证明只激发同相模态：
\end{enumerate}

\[
X_a-X_{a0}=X_b-X_{b0}
=\frac{v_0}{\omega_1}\sin\omega_1t.
\]

若初速度分别为 \(v_0,-v_0\)，则只激发反相模态。

\begin{enumerate}
\def\labelenumi{\arabic{enumi}.}
\setcounter{enumi}{1}
\tightlist
\item
  若上下外侧弹簧劲度分别为 \(k_1,k_2\)，关于平衡位置的刚度矩阵改为
\end{enumerate}

\[
\begin{pmatrix}
k_1+K&-K\\
-K&k_2+K
\end{pmatrix}.
\]

要求写出

\[
\det(\boldsymbol K-m\omega^2\boldsymbol I)=0
\]

并说明 \(k_1\ne k_2\) 时交换对称性破坏，正规模态不再严格为
\((1,\pm1)\)。

\begin{enumerate}
\def\labelenumi{\arabic{enumi}.}
\setcounter{enumi}{2}
\item
  去掉下方弹簧的固定约束、保留上方弹簧。刚度矩阵变为
  \(\bigl(\begin{smallmatrix}k+K&-K\\-K&K\end{smallmatrix}\bigr)\)，行列式为
  \(kK>0\)，故仍无零频整体平移模态。只有两端外弹簧都去掉，才有整体平移的零频方向；竖直自由系统另须处理共同自由落体，不能称静止平衡。
\item
  取 \(K\ll k\)，把
\end{enumerate}

\[
\omega_2=\sqrt{\frac{k+2K}{m}}
\]

展开到 \(K/k\)
的一阶，估计两个频率接近时的拍频。再用给上振子初速度的答案写出振幅包络，联系耦合振子间的能量交换。

\newpage

\section{第 4 题}\label{ux7b2c-4-ux9898}

Gunn 二极管在电压 \(V_1\) 到 \(V_2\) 的工作区间内具有恒定负微分电阻

\[
r_d=\frac{dU}{dI}=-r,\qquad r>0.
\]

射频扼流器 RFC 对直流视为导线、对高频交流视为断路。图示交流回路由
\(L,C,R\) 和负微分电阻 \(-r\) 串联构成。

（1）若扰动电流满足 \(i(0)=0,\ \dot i(0)=\alpha\)，求
\(i(t)\)，给出产生增长振荡的条件和振荡频率。

（2）直流工作电压设为

\[
V=\frac{V_1+V_2}{2},
\]

求器件工作时消耗的平均功率。

\begin{figure}[H]
\centering
\includegraphics[width=0.631\linewidth,height=\textheight,keepaspectratio,alt={图4-1　Gunn 二极管的直流偏置与交流振荡回路}]{../figures/question-04.pdf}
\caption*{图4-1　Gunn 二极管的直流偏置与交流振荡回路}
\end{figure}

\subsection{一、题型定位}\label{ux4e00ux9898ux578bux5b9aux4f4d-3}

负微分电阻使交流储能增长。特征根分别给出增长率和振荡频率；稳态振幅与平均功率还需要非线性信息，不能从起振方程直接确定。

\subsection{二、知识点}\label{ux4e8cux77e5ux8bc6ux70b9-3}

\begin{enumerate}
\def\labelenumi{\arabic{enumi}.}
\tightlist
\item
  小信号模型中 Gunn 二极管等效为电阻 \(-r\)。
\item
  交流回路的总阻尼电阻为 \(R-r\)。
\item
  特征根实部决定振幅增长或衰减，虚部决定振荡频率。
\item
  区间外非线性决定是否稳幅及最终波形；区间宽度只提供电压尺度。
\item
  器件总平均功率等于直流功率与交流平均功率之和。
\end{enumerate}

\subsection{三、物理图像}\label{ux4e09ux7269ux7406ux56feux50cf-3}

正电阻 \(R\) 每周期耗散能量，负微分电阻则把直流电源能量送入交流回路。当
\(r>R\) 时净阻尼为负，小扰动增长；振幅触及 \(V_1,V_2\)
后，小信号模型失效；是否形成稳定周期运动还要看区间外的具体伏安关系。

\subsection{四、解题入口}\label{ux56dbux89e3ux9898ux5165ux53e3-3}

对交流小信号写 KVL，并利用 \(\dot q=i\)
消去电容电荷。得到二阶常系数方程后直接分析特征根。第二问把器件端电压分成直流分量
\(V\) 和交流分量 \(u(t)\)。

\subsection{五、解析}\label{ux4e94ux89e3ux6790-3}

\begin{figure}[H]
\centering
\includegraphics[width=0.707\linewidth,height=\textheight,keepaspectratio,alt={图4-2　线性起振波形、增长包络与特征根条件}]{../figures/solution-04.pdf}
\caption*{图4-2　线性起振波形、增长包络与特征根条件}
\end{figure}

图4-2展示负微分电阻和线性起振包络。增长包络仅说明不稳定，有限周期振幅还需要非线性伏安特性；后文的功率近似另列条件。

\needspace{3\baselineskip}
\textbf{（一）交流微分方程}

取回路电流为 \(i(t)\)，电容电荷为 \(q(t)\)。交流 KVL 为

\[
L\frac{di}{dt}+(R-r)i+\frac{q}{C}=0. \tag{1}
\]

又

\[
\frac{dq}{dt}=i. \tag{2}
\]

对式（1）求导：

\[
L\frac{d^2i}{dt^2}
+(R-r)\frac{di}{dt}
+\frac{i}{C}=0. \tag{3}
\]

除以 \(L\)：

\[
\ddot i+\frac{R-r}{L}\dot i+\frac{1}{LC}i=0. \tag{4}
\]

定义

\[
\gamma=\frac{R-r}{2L},
\qquad
\omega_0=\frac1{\sqrt{LC}}. \tag{5}
\]

特征方程

\[
\beta^2+2\gamma\beta+\omega_0^2=0. \tag{6}
\]

特征根

\[
\beta=-\gamma\pm\sqrt{\gamma^2-\omega_0^2}. \tag{7}
\]

\needspace{3\baselineskip}
\textbf{（二）增长振荡解}

若

\[
\gamma^2<\omega_0^2, \tag{8}
\]

定义

\[
\omega=\sqrt{\omega_0^2-\gamma^2}
=\sqrt{
\frac1{LC}-\frac{(R-r)^2}{4L^2}
}. \tag{9}
\]

通解

\[
i(t)=e^{-\gamma t}
\left(A\cos\omega t+B\sin\omega t\right). \tag{10}
\]

由 \(i(0)=0\) 得 \(A=0\)。再由 \(\dot i(0)=\alpha\) 得

\[
B\omega=\alpha. \tag{11}
\]

故

\[
\boxed{
i(t)=\frac{\alpha}{\omega}
e^{(r-R)t/(2L)}
\sin\omega t
}. \tag{12}
\]

振幅增长要求

\[
\frac{r-R}{2L}>0
\quad\Rightarrow\quad r>R. \tag{13}
\]

同时要保持振荡型特征根：

\[
\frac{(r-R)^2}{4L^2}<\frac1{LC}. \tag{14}
\]

因此增长振荡条件为

\[
\boxed{
R<r<R+2\sqrt{\frac LC}
}. \tag{15}
\]

振荡频率

\[
\boxed{
f=\frac{\omega}{2\pi}
=\frac1{2\pi}
\sqrt{
\frac1{LC}-\frac{(R-r)^2}{4L^2}
}
}. \tag{16}
\]

当 \(r=R\) 时为理想无阻尼振荡；当 \(r<R\) 时振幅衰减。

\needspace{3\baselineskip}
\textbf{（三）直流功率}

RFC 对直流导通，电容对直流断路。若电源电压为 \(V_b\)，器件直流电压为

\[
V=\frac{V_1+V_2}{2}, \tag{17}
\]

则直流电流

\[
I=\frac{V_b-V}{R}. \tag{18}
\]

器件吸收的直流功率

\[
\boxed{
P_D=VI
=\frac{V_1+V_2}{2}
\frac{V_b-\dfrac{V_1+V_2}{2}}{R}
}. \tag{19}
\]

\needspace{3\baselineskip}
\textbf{（四）平均功率能够确定到哪一步}

令实际器件电流为 \(I+j(t)\)，器件电压为
\(V+u(t)\)，交流分量的周期均值均为零。按吸收功率为正，恒有

\[
\overline P=VI+\langle u j\rangle.
\]

若波形始终在恒斜率区内，\(u=-rj\)，则

\[
\overline P=\frac{V(V_b-V)}R-\frac{\langle u^2\rangle}{r}.
\]

这是含交流幅度的条件式。题设没有给出区间外的伏安关系或限幅规律，仅设置直流电压在中点，不能唯一确定
\(\langle u^2\rangle\)。

参考解另取``近似正弦，电压恰好摆满负阻区间''的假设，即
\(u=A\cos\Omega t\)、\(A=(V_2-V_1)/2\)。在此附加近似下才有

\[
\overline P_{\rm ref}=\frac{V_1+V_2}{2R}\left[V_b-\frac{V_1+V_2}{2}\right]
-\frac{(V_2-V_1)^2}{8r}.
\]

不能把该式称为题设唯一的精确稳态功率。对于
\(r>R\)，若一个非零周期波形始终服从
\(u=-rj\)，对交流储能方程取一周期平均，必有

\[
0=(r-R)\langle j^2\rangle>0,
\]

出现矛盾。稳幅必须有区间外非线性、其他耗散或其他加载参与；这些机制会同时改变波形和功率。图4-2的增长曲线仅描绘小信号起振阶段。

\needspace{3\baselineskip}
\textbf{其他特征根分支与初值}

当 \(\gamma^2>\omega_0^2\)，令
\(s=\sqrt{\gamma^2-\omega_0^2}\)，满足同一初值的解为
\(i=(\alpha/s)e^{-\gamma t}\sinh st\)；临界重根时为
\(i=\alpha t e^{-\gamma t}\)。因此 \(r=R\)
是等幅振荡边界，\(r=R+2\sqrt{L/C}\)
是非振荡增长的重根边界。还可由未求导的KVL检查
\(q(0)=-LC\alpha\)；求导后的二阶方程并未允许任意独立选择电容初始扰动。

\needspace{3\baselineskip}
\textbf{用一个非线性例子检验稳幅条件}

为说明缺失信息的作用，可以补设交流器件电压 \(u=-rj+\eta j^3\)，其中
\(\eta>0\)，并近似取 \(j=J\cos\omega t\)。交流储能在一周期后不变，要求

\[
0=(r-R)\langle j^2\rangle-\eta\langle j^4\rangle.
\]

由
\(\langle j^2\rangle=J^2/2\)、\(\langle j^4\rangle=3J^4/8\)，得非零近似幅值

\[
J^2=\frac{4(r-R)}{3\eta}.
\]

同一个线性负阻 \(r\)，只要改变区间外的非线性系数
\(\eta\)，稳态幅值就会改变。因此原题给出的区间端点和斜率，不能唯一代替完整非线性关系。该例子只是补充模型，未把三次伏安关系擅加为原题条件。

如果仍将电压波形按正弦并假设恰好覆盖负阻区间，平均器件功率可按参考式计算；但这项近似和上面的近正弦电流例子应分开使用。对任何正式稳态答案，还需核对每周期负阻供能、正电阻耗能和器件非线性耗能的收支。

\subsection{六、关键洞察}\label{ux516dux5173ux952eux6d1eux5bdf-3}

解析图展示负微分电阻和线性起振包络。增长包络仅说明不稳定，有限周期振幅还需要非线性伏安特性；后文的功率近似另列条件。

\subsection{七、易错点}\label{ux4e03ux6613ux9519ux70b9-3}

\begin{enumerate}
\def\labelenumi{\arabic{enumi}.}
\tightlist
\item
  把负微分电阻当成普通负电源，或把直流伏安关系写成 \(V=-rI\)。\(-r\)
  只适用于工作点附近的交流增量 \(u=-ri\)。
\item
  KVL 中把总电阻写成 \(R+r\)。交流压降为 \(Ri+u=Ri-ri=(R-r)i\)。
\item
  从含 \(q\) 的一阶 KVL到 \(i\) 的二阶方程时漏用
  \(\dot q=i\)，或把电容项写成 \(\dot i/C\)。
\item
  只写 \(r>R\)，漏掉复根条件。这样会把单调指数发散也误称为振荡。
\item
  对不等式开方时漏掉绝对值。一般条件是 \(|R-r|<2\sqrt{L/C}\)，再与
  \(r>R\) 合并才得到式（15）。
\item
  把角频率 \(\omega\) 直接作为频率 \(f\)，漏掉 \(2\pi\)。
\item
  由 \(i(0)=0\) 得到 \(A=0\) 后，求导时漏掉指数包络的项。因为
  \(A=0\)，在 \(t=0\) 处包络项恰不贡献，才有 \(B\omega=\alpha\)。
\item
  把区间宽度直接当成必然的稳态振幅。中点到边界的距离是半宽，但波形是否触及边界、如何越出边界，需要额外条件。
\item
  计算交流平均功率时漏掉 \(\langle\cos^2\omega t\rangle=1/2\)，或因
  \(u=-ri\) 丢失负号。
\item
  把 \(P_A<0\)
  解释为器件``不消耗任何能量''。正确含义仅是交流端口输出能量，总平均功率还要加直流吸收项。
\end{enumerate}

\subsection{八、课堂讲解建议}\label{ux516bux8bfeux5802ux8bb2ux89e3ux5efaux8bae-3}

讲解时先把``起振''和``稳幅''明确分开。

\begin{enumerate}
\def\labelenumi{\arabic{enumi}.}
\tightlist
\item
  起振阶段只用小信号模型。板书微分方程后，把特征根画在复平面上：横坐标是增长率
  \((r-R)/(2L)\)，纵坐标是振荡角频率。让学生分别回答``向右移动意味着什么''\,``虚部消失意味着什么''。
\item
  画参数轴，把 \(r=R\)
  标成无阻尼振荡；右侧先是增长振荡，再是非振荡增长。左侧稳定分支还要按
  \(|R-r|\) 区分衰减振荡与非振荡衰减。这样可以防止只背一个不等式。
\item
  用能量法交叉解释。将 KVL 乘 \(i\)
  后，电感能和电容能组成总储能，\((R-r)i^2\)
  是净耗散率。特征根法给定量解，能量法给方向直觉。
\item
  稳幅阶段在伏安曲线上标出 \(V_1,V_2\) 和中点
  \(V\)。先说明边界仅给电压尺度；非线性细节未给出，参考功率式依赖``正弦摆满区间''的附加近似。再用周期能量平衡检验此近似的局限。
\item
  功率板书固定采用``吸收为正''。先算 \(P_D>0\)，再算
  \(P_A<0\)，最后相加。配合图4-2画能量箭头：偏置电源 \(\to\)
  器件直流端，器件交流端 \(\to\) RLC 回路。
\end{enumerate}

课堂核验可让学生取 \(r=R\)。此时指数包络应消失，频率回到
\(1/(2\pi\sqrt{LC})\)；平均交流功率仍为负阻向回路输出，但在线性边界上恰与
\(R\) 的耗散抵消。

\subsection{九、变式训练}\label{ux4e5dux53d8ux5f0fux8badux7ec3-3}

\begin{enumerate}
\def\labelenumi{\arabic{enumi}.}
\tightlist
\item
  若负微分电阻随电流振幅近似为
\end{enumerate}

\[
r_{\rm eff}(I_m)=r-\eta I_m^2,\qquad \eta>0,
\]

则净阻尼在

\[
r_{\rm eff}(I_m)=R
\]

时为零。由此估计稳定电流振幅

\[
I_m\simeq\sqrt{\frac{r-R}{\eta}}.
\]

说明小振幅时增长、大振幅时衰减如何共同形成稳定极限环。

\begin{enumerate}
\def\labelenumi{\arabic{enumi}.}
\setcounter{enumi}{1}
\item
  若电阻 \(R\) 改为可调，固定 \(L,C,r\)，求起振的 \(R\)
  范围，并讨论在增长振荡区内增大 \(R\) 对增长率和振荡频率的影响。
\item
  给定初值
  \(i(0)=i_0,\ \dot i(0)=0\)，在增长振荡条件下重新确定通解中的两个常数。要求显式保留指数包络求导所产生的项，而不是把答案简单改成余弦。
\item
  若把 RFC
  误当作对交流短路，重画交流等效电路并说明为什么谐振回路会被偏置电源低内阻严重加载。该变式用于巩固``同一器件在直流与交流频段有不同等效连接''的观念。
\end{enumerate}

\newpage

\section{第 5 题}\label{ux7b2c-5-ux9898}

质量为 \(m\)、半径为 \(R\)
的均匀圆盘放在光滑水平桌面上；小虫从盘心沿盘上曲线爬到边缘。原卷未给小虫质量；参考解补取小虫质量为
\(m\)，下面先按此补充条件作答，再给出一般质量结果。以圆盘自身的 \(O\)
点为极点，曲线为

\[
r=\frac{2R}{\pi}\theta,
\qquad 0\le\theta\le\frac{\pi}{2}.
\]

（1）求圆盘相对桌面转过的角度。

（2）求小虫相对桌面、以系统质心为极点的轨道方程。

（3）写出小虫相对桌面路程的积分表达式。

\begin{figure}[H]
\centering
\includegraphics[width=0.618\linewidth,height=\textheight,keepaspectratio,alt={图5-1　小虫在可自由平动转动圆盘上的螺线轨迹}]{../figures/question-05.pdf}
\caption*{图5-1　小虫在可自由平动转动圆盘上的螺线轨迹}
\end{figure}

\subsection{一、题型定位}\label{ux4e00ux9898ux578bux5b9aux4f4d-4}

圆盘会平动，也会转动。先固定系统质心，再合计圆盘自转和两者的轨道角动量；原卷漏写小虫质量，等质量结果须注明这个条件。

\subsection{二、知识点}\label{ux4e8cux77e5ux8bc6ux70b9-4}

\begin{enumerate}
\def\labelenumi{\arabic{enumi}.}
\tightlist
\item
  水平面无外力，系统质心静止。
\item
  圆盘平动与小虫相对运动共同贡献关于系统质心的角动量。
\item
  圆盘转角与小虫相对盘的极角不是同一角。
\item
  极坐标弧长元 \(ds=\sqrt{d\rho^2+\rho^2d\varphi^2}\)。
\end{enumerate}

\subsection{三、物理图像}\label{ux4e09ux7269ux7406ux56feux50cf-4}

小虫向外爬时，圆盘质心向相反方向移动，圆盘也反向旋转以保持总角动量为零。桌面系中小虫到系统质心的距离只有它到盘心距离的一半，因此轨迹不是盘面上原来的阿基米德螺线。

\subsection{四、解题入口}\label{ux56dbux89e3ux9898ux5165ux53e3-4}

令圆盘转角为 \(\phi\)，小虫相对圆盘极角为
\(\theta\)，则小虫在桌面系的极角为
\(\varphi=\phi+\theta\)。由两者质量相等，系统质心到小虫的距离为
\(\rho=r/2\)。

\subsection{五、解析}\label{ux4e94ux89e3ux6790-4}

\begin{figure}[H]
\centering
\includegraphics[width=0.86\linewidth,height=\textheight,keepaspectratio,alt={图5-2　系统质心、盘心和小虫位置矢量的关系}]{../figures/solution-05.pdf}
\caption*{图5-2　系统质心、盘心和小虫位置矢量的关系}
\end{figure}

图5-2按 \(m_b=m\)
的参考假设画虫相对总质心的轨迹，末态盘心在相反方向。图中径向距离和方位角分别对应
\(\rho\)、\(\psi=\theta+\phi\)。

\needspace{3\baselineskip}
\textbf{（一）关于系统质心的角动量}

由于圆盘与小虫质量都为 \(m\)，若 \(\boldsymbol r\)
是小虫相对盘心的位置，则相对系统质心的位置分别为

\[
\boldsymbol r_{\rm bug}=\frac{\boldsymbol r}{2},
\qquad
\boldsymbol r_{\rm disk}=-\frac{\boldsymbol r}{2}. \tag{1}
\]

两者质心平动产生的总轨道角动量为

\[
L_{\rm orb}
=\frac{m}{2}r^2\dot\varphi. \tag{2}
\]

圆盘绕自身中心的转动惯量

\[
I_d=\frac12mR^2. \tag{3}
\]

其自转角动量为

\[
L_{\rm spin}=\frac12mR^2\dot\phi. \tag{4}
\]

初态静止，且无外力矩，所以

\[
\frac{m}{2}r^2\dot\varphi
+\frac12mR^2\dot\phi=0. \tag{5}
\]

又

\[
\varphi=\phi+\theta. \tag{6}
\]

代入并约去 \(m/2\)：

\[
r^2(\dot\phi+\dot\theta)+R^2\dot\phi=0. \tag{7}
\]

所以

\[
\boxed{
\dot\phi=-\frac{r^2}{r^2+R^2}\dot\theta
}. \tag{8}
\]

\needspace{3\baselineskip}
\textbf{（二）圆盘总转角}

代入 \(r=2R\theta/\pi\)：

\[
\phi_{\rm f}
=-\int_0^{\pi/2}
\frac{(2\theta/\pi)^2}
{1+(2\theta/\pi)^2}\,d\theta. \tag{9}
\]

令

\[
u=\frac{2\theta}{\pi},
\qquad d\theta=\frac{\pi}{2}du. \tag{10}
\]

则

\[
\begin{aligned}
\phi_{\rm f}
&=-\frac{\pi}{2}
\int_0^1\frac{u^2}{1+u^2}\,du\\
&=-\frac{\pi}{2}
\int_0^1\left(1-\frac1{1+u^2}\right)du\\
&=-\frac{\pi}{2}
\left[u-\arctan u\right]_0^1\\
&=-\frac{\pi}{2}\left(1-\frac{\pi}{4}\right).
\end{aligned} \tag{11}
\]

因此

\[
\boxed{
\phi_{\rm f}
=-\frac{\pi}{2}\left(1-\frac{\pi}{4}\right)
}. \tag{12}
\]

负号表示圆盘转动方向与小虫相对爬行角方向相反。

\needspace{3\baselineskip}
\textbf{（三）桌面系轨道方程}

小虫到系统质心的极径

\[
\rho=\frac r2=\frac{R}{\pi}\theta. \tag{13}
\]

由式（6）、式（8）

\[
\dot\varphi
=\dot\theta+\dot\phi
=\frac{R^2}{r^2+R^2}\dot\theta. \tag{14}
\]

所以

\[
\frac{d\varphi}{d\theta}
=\frac1{1+(2\theta/\pi)^2}. \tag{15}
\]

积分：

\[
\begin{aligned}
\varphi
&=\int_0^\theta
\frac{d\theta'}{1+(2\theta'/\pi)^2}\\
&=\frac{\pi}{2}
\arctan\frac{2\theta}{\pi}.
\end{aligned} \tag{16}
\]

而 \(\theta=\pi\rho/R\)，故

\[
\varphi=\frac{\pi}{2}
\arctan\frac{2\rho}{R}. \tag{17}
\]

反解：

\[
\boxed{
\rho=\frac R2\tan\frac{2\varphi}{\pi},
\qquad 0\le\varphi\le\frac{\pi^2}{8}
}. \tag{18}
\]

\needspace{3\baselineskip}
\textbf{（四）路程积分}

极坐标弧长元

\[
ds=\sqrt{d\rho^2+\rho^2d\varphi^2}. \tag{19}
\]

由式（13）

\[
\frac{d\rho}{d\theta}=\frac R\pi. \tag{20}
\]

由式（15）

\[
\frac{d\varphi}{d\theta}
=\frac1{1+4\theta^2/\pi^2}. \tag{21}
\]

所以

\[
\boxed{
s=
\frac R\pi
\int_0^{\pi/2}
\sqrt{
1+
\left[
\frac{\theta}
{1+4\theta^2/\pi^2}
\right]^2
}\,d\theta
}. \tag{22}
\]

题目只要求积分表达式，无需强行求初等原函数。

\needspace{3\baselineskip}
\textbf{补充推导与适用条件}

式（2）的系数 \(m/2\)
可以用约化质量理解。对任意两个质点的相对运动，关于系统质心的轨道角动量是

\[
L_{\rm orb}=\mu r^2\dot\varphi,
\qquad
\mu=\frac{m\cdot m}{m+m}=\frac m2.
\]

圆盘除质心轨道运动外，还有绕自身中心的自转角动量
\(I_d\dot\phi\)。这两部分都必须计算；只写
\(I_d\dot\phi+mr^2\dot\varphi=0\) 会同时错在参考点和质量系数。

由

\[
\frac{d\phi}{d\theta}
=-\frac{r^2}{r^2+R^2}
\]

可看出圆盘转角只取决于小虫在盘上的路径
\(r(\theta)\)，与小虫爬行快慢无关。这是一种几何型结果：只要小虫可视为质点且无水平外力及外力矩，沿同一条曲线以不同时间规律爬行，最终圆盘转角相同。

终点可作几何一致性检查。由式（12）

\[
\phi_f
=-\frac{\pi}{2}\left(1-\frac{\pi}{4}\right),
\]

而小虫相对盘转过 \(\theta_f=\pi/2\)，所以桌面系极角

\[
\varphi_f=\theta_f+\phi_f=\frac{\pi^2}{8},
\]

恰好等于轨道方程式（18）给出的上限。

路程必须沿桌面系轨迹计算。圆盘本身平移、转动，使桌面系轨迹与盘面螺线不同。式（22）中的第一项来自极径变化，第二项来自方位角变化；两项都是长度平方，开方后乘
\(d\theta\) 才得到实际弧长。

还可检查数量级。因为

\[
0\le
\frac{\theta}{1+4\theta^2/\pi^2}
\le\frac{\pi}{4},
\]

所以被积函数介于 \(1\) 与 \(\sqrt{1+\pi^2/16}\)
之间，路程应略大于径向位移 \(R/2\)，不会大到多个圆周。

\needspace{3\baselineskip}
\textbf{原卷缺项与一般质量答案}

设小虫质量为 \(m_b\)，定义
\(\beta=m/(m+m_b)\)、\(\mu=mm_b/(m+m_b)\)，以及
\(\eta=2m_b/(m+m_b)\)。图5-2中的位置关系应改为
\(\rho=\beta r\)；等质量时才是 \(r/2\)。守恒式为

\[
\mu r^2\dot\varphi+\tfrac12mR^2\dot\phi=0,
\qquad \varphi=\theta+\phi.
\]

因而沿本题螺线

\[
\varphi(\theta)=\frac{\pi}{2\sqrt\eta}
\arctan\left(\frac{2\sqrt\eta}{\pi}\theta\right),
\qquad \phi(\theta)=\varphi(\theta)-\theta.
\]

三个一般结果为

\[
\phi_f=-\frac\pi2+\frac{\pi}{2\sqrt\eta}\arctan\sqrt\eta,
\]

\[
\rho=\frac{\beta R}{\sqrt\eta}
\tan\left(\frac{2\sqrt\eta}{\pi}\varphi\right),
\quad
0\le\varphi\le\frac{\pi}{2\sqrt\eta}\arctan\sqrt\eta,
\]

\[
s=\frac{2\beta R}{\pi}\int_0^{\pi/2}
\sqrt{1+\left[\frac{\theta}{1+4\eta\theta^2/\pi^2}\right]^2}\,d\theta.
\]

令 \(m_b=m\)，则 \(\beta=1/2,\eta=1\)，逐项恢复主答案。令
\(m_b/m\to0\)，盘转角趋零，小虫轨迹趋于固定盘面螺线；质量比会改变全部结果，故原题没有补充质量时不能报出唯一数值转角。小虫的爬行速度不必准静态：守恒式本身对任意光滑时间参数化成立，但不提供爬行所需的力或功。

\needspace{3\baselineskip}
\textbf{一般质量比下的初末几何检查}

设虫质量为
\(m_b\)，\(\beta=m/(m+m_b)\)、\(\eta=2m_b/(m+m_b)\)。虫相对盘心的距离是
\(r=2R\theta/\pi\)，相对总质心则为
\(\rho=\beta r\)；盘心相对质心的位置为
\(-(1-\beta)\bm r\)。两者在每一时刻都满足

\[
m\bm r_O+m_b\bm r_b=0.
\]

在等质量参考分支中，\(\beta=1/2,\eta=1\)，虫和盘心距质心相等、方向相反。虫走到盘沿时
\(\theta=\pi/2\)，盘的转角为

\[
\phi_f=-\frac\pi2+\frac{\pi^2}{8},\qquad
\psi_f=\frac{\pi^2}{8}=70.686^\circ.
\]

这正是图5-2末态两点所用的方向。转角应记为弧度式后再换算角度，不能把
\(\pi^2/8\) 当作以度计的数。

当
\(m_b/m\to0\)，\(\eta\to0,\beta\to1\)，盘转角趋于零，虫的质心轨迹趋于原盘心坐标下的螺线。质量比改变后不能保留等质量图中的
\(\rho=r/2\)；路程积分也要同步用一般 \(\rho(\theta),\psi(\theta)\)。

\subsection{六、关键洞察}\label{ux516dux5173ux952eux6d1eux5bdf-4}

解析图按 \(m_b=m\)
的参考假设画虫相对总质心的轨迹，末态盘心在相反方向。图中径向距离和方位角分别对应
\(\rho\)、\(\psi=\theta+\phi\)。

\subsection{七、易错点}\label{ux4e03ux6613ux9519ux70b9-4}

\begin{enumerate}
\def\labelenumi{\arabic{enumi}.}
\tightlist
\item
  把盘心 \(O\)
  当成桌面系固定点。盘心会因小虫爬行而平移；以它为原点写守恒定律会漏掉惯性项。
\item
  由``系统质心静止''错误地推出``圆盘不动''。质心静止只约束盘心和小虫的位置加权和，不排除二者相对运动以及圆盘自转。
\item
  把小虫到系统质心的距离写成 \(r\)。等质量时应为
  \(r/2\)；一般质量比下更不是简单的一半。
\item
  把桌面系极角写成 \(\theta\)。\(\theta\) 是相对盘的极角，盘又转过
  \(\phi\)，所以 \(\varphi=\theta+\phi\)。
\item
  角动量只计圆盘自转，漏掉盘心与小虫绕系统质心的轨道角动量。错误方程通常会给出过大的反向转角。
\item
  在式（8）积分时把 \(r\) 当常数。小虫沿螺线爬行，\(r\) 是 \(\theta\)
  的函数，必须先代入 \(r=2R\theta/\pi\)。
\item
  变量代换 \(u=2\theta/\pi\) 后漏掉
  \(d\theta=(\pi/2)du\)，导致最终转角少一个 \(\pi/2\)。
\item
  轨道方程只给参数式 \(\rho(\theta),\varphi(\theta)\) 而未按要求消去
  \(\theta\)。参数式可作为中间结果，最终还应得到式（18）。
\item
  用盘面曲线弧长
\end{enumerate}

\[
\int\sqrt{dr^2+r^2d\theta^2}
\]

回答第三问。题目要求小虫相对桌面的路程，必须使用 \(\rho,\varphi\)。

\subsection{八、课堂讲解建议}\label{ux516bux8bfeux5802ux8bb2ux89e3ux5efaux8bae-4}

建议先做一维类比，再进入极坐标。

\begin{enumerate}
\def\labelenumi{\arabic{enumi}.}
\tightlist
\item
  让学生回忆人在光滑小车上向右走时，小车向左移且系统质心不动。把小车换成圆盘后，平动约束完全相同，只是又多了绕盘心的转动。
\item
  在图5-2上固定系统质心，画出盘心与小虫的相反位置矢量。先写
\end{enumerate}

\[
m\boldsymbol r_{\rm bug}+m\boldsymbol r_{\rm disk}=0,
\]

再结合两者之差为 \(\boldsymbol r\)，解出 \(\pm\boldsymbol r/2\)。 3.
分两行写角动量：第一行是二者质心轨道角动量，第二行是圆盘自转角动量。让学生逐项说出``质量、力臂平方、角速度''，避免直接接受式（5）。
4. 明确三个角的含义：\(\theta\) 随小虫相对盘爬行而变，\(\phi\)
是盘相对桌面转角，\(\varphi\) 是小虫相对桌面的方位角。把关系
\(\varphi=\phi+\theta\) 写在图旁全程保留。 5. 求轨道时先得到参数式

\[
\rho=\frac{R\theta}{\pi},\qquad
\varphi=\frac{\pi}{2}\arctan\frac{2\theta}{\pi},
\]

再消元。这样学生能看清轨道方程来自哪两条关系。 6.
弧长部分从小位移直角三角形出发解释
\(ds^2=d\rho^2+\rho^2d\varphi^2\)，而不是直接背极坐标公式。

结尾用端点复核：\(\theta=0\) 时 \(\rho=\varphi=0\)；\(\theta=\pi/2\) 时
\(\rho=R/2,\ \varphi=\pi^2/8\)。若任何一项不符，应回查角度合成或质心距离。

\subsection{九、变式训练}\label{ux4e5dux53d8ux5f0fux8badux7ec3-4}

\begin{enumerate}
\def\labelenumi{\arabic{enumi}.}
\tightlist
\item
  若圆盘质量为 \(M_d\)、小虫质量为 \(m_b\)，系统质心到小虫的距离为
\end{enumerate}

\[
\rho=\frac{M_d}{M_d+m_b}r,
\]

相对运动的约化质量为

\[
\mu=\frac{M_dm_b}{M_d+m_b}.
\]

证明角动量守恒给出

\[
\dot\phi
=-\frac{\mu r^2}{I_d+\mu r^2}\dot\theta,
\qquad
I_d=\frac12M_dR^2.
\]

再令 \(M_d=m_b=m\)，验证回到式（8）。

\begin{enumerate}
\def\labelenumi{\arabic{enumi}.}
\setcounter{enumi}{1}
\tightlist
\item
  若小虫沿圆周 \(r=r_0\) 爬过相对角
  \(\Theta\)，求圆盘转角和小虫桌面系极角。由于系数为常数，应有
\end{enumerate}

\[
\phi
=-\frac{r_0^2}{r_0^2+R^2}\Theta.
\]

讨论 \(r_0\ll R\) 与 \(r_0=R\) 两个极限。

\begin{enumerate}
\def\labelenumi{\arabic{enumi}.}
\setcounter{enumi}{2}
\tightlist
\item
  若圆盘与桌面之间存在很小的恒定阻力矩
  \(\tau_f\)，说明总角动量不再守恒，并写出
\end{enumerate}

\[
\frac{dL}{dt}=-\tau_f\,\operatorname{sgn}\dot\phi.
\]

此时最终转角为何会依赖小虫爬行所用时间？该变式用于辨认本题``路径无关速度''的成立条件。

\begin{enumerate}
\def\labelenumi{\arabic{enumi}.}
\setcounter{enumi}{3}
\tightlist
\item
  对式（22）作数值积分，得到路程的近似值；再分别与盘面螺线长度和直线距离
  \(R/2\) 比较。数值程序应同时检查被积函数在两端连续且为正。
\end{enumerate}

\newpage

\section{第 6 题}\label{ux7b2c-6-ux9898}

黑盒有三个接线柱，内部含电容 \(C\)、带内阻 \(r\) 的电感 \(L\) 和电阻
\(R=100\,\Omega\)。已知

\[
1.0\,\mu{\rm F}\le C\le10\,\mu{\rm F},\quad
10\,{\rm mH}\le L\le30\,{\rm mH},
\]

\[
10\,\Omega\le r\le30\,\Omega.
\]

按题图接入信号源与示波器。扫频时 \(CH_1\) 幅度随频率增大而增大；在
\(1\,\mathrm{kHz}\) 下记录两幅李萨如图，第二幅由交换黑盒节点 1、2 得到。

（1）判断内部连接。

（2）由图中数据求 \(L,r,C\)，保留三位有效数字。

（3）求 \(LC\) 串联谐振频率。

\begin{figure}[H]
\centering
\includegraphics[width=0.9\linewidth,height=\textheight,keepaspectratio,alt={图6-1　黑盒测试电路、四种候选拓扑和两幅李萨如图}]{../figures/question-06.pdf}
\caption*{图6-1　黑盒测试电路、四种候选拓扑和两幅李萨如图}
\end{figure}

\subsection{一、题型定位}\label{ux4e00ux9898ux578bux5b9aux4f4d-5}

高阻测量端使一条支路几乎无电流，交换端子后分别得到RC和RL分压。椭圆给出幅值与相位大小；相位的正负要由电路或扫描方向判断。

\subsection{二、知识点}\label{ux4e8cux77e5ux8bc6ux70b9-5}

\begin{enumerate}
\def\labelenumi{\arabic{enumi}.}
\tightlist
\item
  扫频幅频趋势可区分电容支路和电感支路。
\item
  两同频正弦信号形成椭圆，截距比给出相位差。
\item
  交换节点使同一三端网络分别呈现 \(RC\) 和 \(RL\) 测量状态。
\item
  复数阻抗能统一处理幅值比与相位差。
\item
  串联谐振频率 \(f_0=1/(2\pi\sqrt{LC})\)。
\end{enumerate}

\subsection{三、物理图像}\label{ux4e09ux7269ux7406ux56feux50cf-5}

选项 B 的中间节点分别连接 \(L+r\)、\(R\) 和
\(C\)。一种接法使电感支路近似开路，测得 \(RC\)
相位；交换节点后使电容支路近似开路，测得 \(RL\)
相位。两次测量共同确定三个参数。

\subsection{四、解题入口}\label{ux56dbux89e3ux9898ux5165ux53e3-5}

对李萨如图设

\[
x=A\cos\omega t,\qquad
y=B\cos(\omega t+\varphi).
\]

利用横、纵轴截距求 \(\sin\varphi\)，再由幅值标尺求
\(B/A\)。把实验比值与相应分压复数比相等。

\subsection{五、解析}\label{ux4e94ux89e3ux6790-5}

\begin{figure}[H]
\centering
\includegraphics[width=0.9\linewidth,height=\textheight,keepaspectratio,alt={图6-2　两次接法的等效分压电路与复阻抗三角形}]{../figures/solution-06.pdf}
\caption*{图6-2　两次接法的等效分压电路与复阻抗三角形}
\end{figure}

图6-2中两次等效接法分别测RC、RL分压，阻抗三角形的边为
\(R+r,\omega L,|R/H_2|\)。读数反演用同一组椭圆截距，避免中途舍入相位。

\needspace{3\baselineskip}
\textbf{（一）拓扑判断}

测量电路中，3 端接公共地，2 端接信号源，\(CH_1\) 测量 1
端对地电压，\(CH_2\) 测量信号源电压。示波器输入阻抗很大，可以把 1
端看成近似开路。

若内部为拓扑 B，1 端通过带内阻的电感接到中间节点，中间节点再经 \(R\)
接地、经 \(C\) 接到信号源。因为 1 端开路，电感中几乎没有电流，故 1
端电势等于中间节点电势。第一次接法实际只剩下串联 \(C\)、\(R\)，且
\(CH_1\) 测量 \(R\) 两端电压：

\[
\frac{V_{CH_1}}{V_{CH_2}}
=\frac{R}{R+1/(i\omega C)}
=\frac{i\omega RC}{1+i\omega RC}. \tag{1}
\]

当频率升高时，电容容抗 \(1/(\omega C)\) 减小，式（1）的模从 0 单调趋近
1，恰好符合题述``\(CH_1\)
幅度逐渐增大''的高通特征。候选A的LC通路会出现谐振，幅频响应不能在所扫频段单调增长；C、D在该高阻接法下给出零输出。与非零单调高通现象相容的是B，因此

\[
\boxed{\text{内部连接选 B}}. \tag{2}
\]

\needspace{3\baselineskip}
\textbf{（二）李萨如图的通用读数}

若

\[
x=A\cos\omega t,\qquad
y=B\cos(\omega t+\varphi), \tag{3}
\]

则 \(A\)、\(B\)
分别是横轴信号和纵轴信号的振幅。它们是椭圆在坐标方向上的最大投影，不是椭圆长、短半轴。令
\(y=0\)，有 \(\omega t+\varphi=\pi/2+n\pi\)，从而

\[
|x_0|=A|\sin\varphi|. \tag{4}
\]

同理，令 \(x=0\) 可得

\[
|y_0|=B|\sin\varphi|. \tag{5}
\]

因此

\[
|\sin\varphi|
=\frac{|x_0|}{A}
=\frac{|y_0|}{B}. \tag{6}
\]

椭圆倾斜方向决定 \(\cos\varphi\) 的符号，静态椭圆不能区分 \(+\varphi\)
和
\(-\varphi\)；超前或滞后要从电路性质或标出的时间扫描方向判断。横、纵轴标尺给出
\(A,B\)。同一幅图可用两个截距各算一次
\(|\sin\varphi|\)，二者应在曲线线宽允许的误差内一致，这也是读图的自检方法。

\needspace{3\baselineskip}
\textbf{（三）第一次接法：求电容}

图 1 中横轴为 \(CH_1\)、纵轴为 \(CH_2\)。由图读出

\[
A_1=3.54\,\mathrm V,\qquad B_1\simeq5.00\,\mathrm V,
\]

\[
x_{01}=2.50\,\mathrm V,\qquad y_{01}=3.53\,\mathrm V.
\]

于是

\[
\frac{x_{01}}{A_1}
=\frac{2.50}{3.54}
\simeq0.706,
\]

\[
\frac{y_{01}}{B_1}
=\frac{3.53}{5.00}
\simeq0.706.
\]

两种读法相符，故

\[
|\varphi_1|\simeq\arcsin(0.706)\simeq44.9^\circ. \tag{7}
\]

幅值比为

\[
|H_1|=\frac{A_1}{B_1}
=\frac{3.54}{5.00}
\simeq0.708. \tag{8}
\]

对式（1）令 \(x=\omega RC\)，有

\[
|H_1|=\frac{x}{\sqrt{1+x^2}},
\qquad
\arg H_1=\arctan\frac1x.
\]

由相位和幅值两种读法都得到

\[
\omega RC\simeq1.00. \tag{9}
\]

\(f=1.000\,\mathrm{kHz}\)，故

\[
\begin{aligned}
C
&=\frac{1}{2\pi fR}\\
&=\frac{1}{2\pi(1000)(100)}\\
&\simeq1.59\times10^{-6}\,\mathrm F.
\end{aligned} \tag{10}
\]

结合图上有效数字，

\[
\boxed{C=1.60\,\mu\mathrm F}. \tag{11}
\]

\needspace{3\baselineskip}
\textbf{（四）交换节点：求电感和内阻}

把黑盒的 1、2 端交换后接回同一测量电路。此时信号源从原来的 1 端经
\(r+i\omega L\) 接到中间节点，\(CH_1\) 从原来的 2
端测量中间节点；由于示波器近似开路，电容支路没有电流。故第二次接法是
\(r+i\omega L\) 与 \(R\) 串联、输出取在 \(R\) 上：

\[
H_2
=\frac{V_{CH_1}}{V_{CH_2}}
=\frac{R}{R+r+i\omega L}. \tag{12}
\]

注意：分子是 \(R\)，不是 \(r+i\omega L\)；\(CH_1\)
测的是中间节点对地电压，也就是电阻电压。

由图 2 读出

\[
A_2=3.95\,\mathrm V,\qquad B_2\simeq5.00\,\mathrm V,
\]

以及

\[
x_{02}=1.96\,\mathrm V,\qquad y_{02}=2.48\,\mathrm V.
\]

因此

\[
\frac{x_{02}}{A_2}
=\frac{1.96}{3.95}
\simeq0.496,
\qquad
\frac{y_{02}}{B_2}
=\frac{2.48}{5.00}
\simeq0.496. \tag{13}
\]

令输出相对输入的滞后角大小为 \(\alpha\)，则

\[
\alpha=\arcsin(1.96/3.95)\simeq29.749^\circ. \tag{14}
\]

幅值比

\[
|H_2|=\frac{A_2}{B_2}
=\frac{3.95}{5.00}
\simeq0.790. \tag{15}
\]

将式（12）写成极坐标形式：

\[
H_2
=\frac{R}{\sqrt{(R+r)^2+(\omega L)^2}}
\exp(-i\alpha).
\]

因而

\[
\frac{R}{H_2}
=(R+r)+i\omega L
\]

在复平面上的模和辐角与分母相同。分别投影到实轴、虚轴：

\[
R+r=\frac{R}{|H_2|}\cos\alpha, \tag{16}
\]

\[
\omega L=\frac{R}{|H_2|}\sin\alpha. \tag{17}
\]

于是

\[
\begin{aligned}
r
&=100\frac{2.48}{1.96}\sqrt{1-(1.96/3.95)^2}-100\\
&\simeq9.85\,\Omega,
\end{aligned} \tag{18}
\]

\[
\begin{aligned}
L&=\frac{100(2.48)}{(2\pi1000)(3.95)}\\
&=9.99251\times10^{-3}\,\mathrm H\simeq9.99\,\mathrm{mH}.
\end{aligned}\tag{19}
\]

图线本身有宽度，\(A_2,B_2,x_{02},y_{02}\)
的末位不是精确测量值；在图线允许的读数区间内，\(L\) 约为
\(10.0\,\mathrm{mH}\)。按下面不预先舍入相位的反演中心值，保留三位有效数字：

\[
\boxed{
L=9.99\,\mathrm{mH},\quad
r=9.85\,\Omega,\quad
C=1.60\,\mu\mathrm F
}. \tag{20}
\]

其中 \(r\) 的中心读数略低于题给范围下限 \(10\,\Omega\)，偏差约
\(1.5\%\)；不能把题给范围与反演中心值同时当作严格精确数据。需报告读数不确定度，或补查原始测量数据，不应把中心值强行截成
\(10.0\,\Omega\)。

\needspace{3\baselineskip}
\textbf{（五）串联谐振频率}

\[
f_0=\frac1{2\pi\sqrt{LC}}. \tag{21}
\]

代入

\[
L=9.99\times10^{-3}\,\mathrm H,\qquad
C=1.60\times10^{-6}\,\mathrm F:
\]

\[
\begin{aligned}
f_0
&=\frac1{2\pi
\sqrt{(9.99\times10^{-3})(1.60\times10^{-6})}}\\
&\simeq1.26\times10^3\,\mathrm{Hz}.
\end{aligned} \tag{22}
\]

所以

\[
\boxed{f_0=1.26\times10^3\,\mathrm{Hz}}. \tag{23}
\]

\needspace{3\baselineskip}
\textbf{补充推导与适用条件}

\begin{enumerate}
\def\labelenumi{\arabic{enumi}.}
\tightlist
\item
  扫频趋势只负责筛选拓扑。\(CH_1\)
  随频率升高而增大，说明第一次接法呈高通，定位到电容串联、输出取电阻的等效结构。
\item
  第一幅李萨如图给出 \(RC\) 分压的幅值比和相位差，因 \(R\)
  已知，所以确定 \(C\)。
\item
  交换 1、2
  端后，原来开路的电感支路进入输入通道，原来的电容支路变成开路测量端。第二幅图于是确定
  \(r\) 与 \(L\)。
\item
  最后一问只是把已反演的 \(L,C\) 代入独立的串联谐振公式。
\end{enumerate}

示波器输入阻抗很大是等效电路成立的隐含条件。第一次接法中 1 端只有
\(CH_1\)，故电感支路近似无电流，电感两端没有压降；第二次接法中原 2
端只有
\(CH_1\)，故电容支路近似无电流。所谓``交换节点''，实质是利用高阻测量端选择性地让某一支路退出电流回路。

李萨如图同时包含两类独立信息：

\[
\frac AB=|H|
\]

给传递函数的模，

\[
\frac{x_0}{A}=\frac{y_0}{B}=\sin|\varphi|
\]

给相位差。只用截距比只能得到相位，只用最大投影比只能得到幅值；求 \(r,L\)
需要二者同时参与。

第二次接法的复阻抗三角形是最简洁的反演图像。由

\[
\frac{R}{H_2}=R+r+i\omega L
\]

可知斜边为 \(R/|H_2|\)，水平边为 \(R+r\)，竖直边为
\(\omega L\)。实部与虚部投影直接给式（16）、式（17），比把复分式有理化后联立两条长方程更清楚。

图读数必须承认不确定度。标注中心值 \(r=9.85\,\Omega\)
低于题给下界；图线宽度或末位量化可能解释差异，但没有原始测量数据时不能确定原因。应报告中心值及明确假定的读数区间，保留数据之间的差异。

数值还可做反代检查。取 \(r=9.85\,\Omega,L=9.99\,\mathrm{mH}\)，在
\(1\,\mathrm{kHz}\) 时

\[
R+r\simeq109.8\,\Omega,\qquad
\omega L\simeq62.8\,\Omega,
\]

所以相位约为 \(\arctan(62.8/109.8)=29.8^\circ\)，幅值比约为
\(100/\sqrt{109.8^2+62.8^2}=0.790\)，与图 2 同时吻合。

\needspace{3\baselineskip}
\textbf{直接用标注读数反演，避免中途舍入}

图1的纵轴最大投影没有单独标出
\(5.00\,\mathrm V\)。在RC模型下，\(\sin\alpha_1=2.50/3.54\)，且
\(|H_1|=\cos\alpha_1\)；由此推得输入幅度约
\(5.00\,\mathrm V\)，不是凭图形高度任意增添的测量量。

\[
C=\frac1{\omega R\tan\alpha_1}=1.59557\,\mu\mathrm F.
\]

图2以标注 \(A=3.95\,\mathrm V,x_0=1.96\,\mathrm V,y_0=2.48\,\mathrm V\)
计算：\(\sin\alpha=x_0/A\)，\(B=Ay_0/x_0=4.99796\,\mathrm V\)。不把
\(\alpha\) 先舍入为 \(29.8^\circ\)，则

\[
r=R\frac{y_0}{x_0}\sqrt{1-\frac{x_0^2}{A^2}}-R=9.85474\,\Omega,
\]

\[
L=\frac{R y_0}{\omega A}=9.99251\,\mathrm{mH},\qquad
f_0=1260.45\,\mathrm{Hz}.
\]

所以中心值为
\(r=9.85\,\Omega,L=9.99\,\mathrm{mH},C=1.60\,\mu\mathrm F\)。参考解中的约
\(9.80\)---\(9.85\,\Omega\)
是读图答案范围，不能用它代替原始标注的逐项计算。

若仅把三个两位小数读数各看作四舍五入到 \(0.01\,\mathrm V\)，在其
\(\pm0.005\,\mathrm V\) 区间的八个角点计算，所得 \(r\) 范围约为
\(9.22\)---\(10.49\,\Omega\)。这是明确假定末位量化误差后的区间，足以覆盖
\(10\,\Omega\)；未给仪器与线宽数据时不能称为实测置信区间。图6-1的曲线与截距应由同一组读数生成，不能画任意倾斜椭圆再标上这些数字。

\subsection{六、关键洞察}\label{ux516dux5173ux952eux6d1eux5bdf-5}

解析图中两次等效接法分别测RC、RL分压，阻抗三角形的边为
\(R+r,\omega L,|R/H_2|\)。读数反演用同一组椭圆截距，避免中途舍入相位。

\subsection{七、易错点}\label{ux4e03ux6613ux9519ux70b9-5}

\begin{enumerate}
\def\labelenumi{\arabic{enumi}.}
\tightlist
\item
  没有先标清外部接线，直接从候选图猜 B。应明确 3 端接地、2 端输入、1
  端高阻测量，才能判断哪条支路开路。
\item
  认为``开路支路不存在''。开路支路中电流近似为零，但其端点电势仍能被示波器测量；正因为没有压降，测量端电势等于相邻节点。
\item
  把第一次接法写成低通
\end{enumerate}

\[
\frac1{1+i\omega RC}.
\]

\(CH_1\) 取的是 \(R\) 上电压，正确传递函数是高通的
\(i\omega RC/(1+i\omega RC)\)。 4. 把椭圆的几何长半轴、短半轴当作
\(A,B\)。\(A,B\) 是椭圆在横、纵坐标方向的最大投影。 5. 把横轴截距
\(x_0\) 直接当作横轴振幅 \(A\)。只有在相位差 \(90^\circ\) 时二者才相等。
6. 只取 \(\arcsin\) 的主值而不看椭圆倾斜方向。截距给
\(|\sin\varphi|\)，相位超前或滞后还要由轨迹方向或电路的容、感性质确定。
7. 第二次接法把分压分子写成 \(r+i\omega L\)。\(CH_1\) 测的是 \(R\)
上电压，故分子必须是
\(R\)；这一错误虽然可能凑出相近数值，却没有正确电路含义。 8. 在求 \(L\)
时使用 \(fL\)，漏掉 \(\omega=2\pi f\)。 9. 把 \(3.95,2.48,1.96\)
三个读数不加区分地代入同一比例。应先标明哪一个是振幅投影、哪两个是坐标轴截距。
10. 计算谐振频率时未把 \(\mathrm{mH}\)、\(\mu\mathrm F\) 转成
SI，导致结果相差 \(10^3\) 或更多。

\subsection{八、课堂讲解建议}\label{ux516bux8bfeux5802ux8bb2ux89e3ux5efaux8bae-5}

先在候选B上描出实际电流路径，再从椭圆的截距和轴向投影求相位与幅值。计算结束后回代两次接法的传递函数。

\begin{enumerate}
\def\labelenumi{\arabic{enumi}.}
\tightlist
\item
  在题图上用三种颜色描出地线、输入线、\(CH_1\) 测量线。先让学生在候选 B
  上沿电流路径走一遍：第一次是输入端 2 经 \(C\)、\(R\)
  到地；电感末端只接高阻示波器，因此不在电流回路中。
\item
  交换端子后重新描线，不允许口头说``显然变成 RL''。明确电流从原 1 端经
  \(L+r\)、\(R\) 到地，而原 2 端的电容只把节点电势送到 \(CH_1\)。
\item
  在标准椭圆上分别标出 \(A,B,x_0,y_0\)。用一组简单相位
  \(\varphi=30^\circ\) 现场验证 \(x_0/A=y_0/B=1/2\)，建立读图公式。
\item
  对图 1 先做两次相位读数：
\end{enumerate}

\[
\frac{2.50}{3.54},\qquad
\frac{3.53}{5.00}.
\]

二者接近说明标注没有认反。再用 \(A_1/B_1\) 检查幅值关系。 5. 对图 2
不直接有理化复分式，而使用图6-2中的复阻抗直角三角形。斜边、邻边、对边分别对应
\(R/|H_2|\)、\(R+r\)、\(\omega L\)，学生只需一次三角投影。 6.
数值求出后立刻反代相位和幅值。若只验证最终谐振频率，无法发现 \(r,L\)
在反演中被交换或分子写错。

可安排两人互检：一人只负责从图上读
\(A,B,x_0,y_0\)，另一人只负责画等效电路和写传递函数，最后交换检查。这样能把``读错图''和``列错电路''两类错误分离。

\subsection{九、变式训练}\label{ux4e5dux53d8ux5f0fux8badux7ec3-5}

\begin{enumerate}
\def\labelenumi{\arabic{enumi}.}
\tightlist
\item
  若示波器两通道增益未知但稳定，将同一正弦信号同时送入两通道，记录显示振幅比
\end{enumerate}

\[
g=\frac{A_{\rm disp}}{B_{\rm disp}}.
\]

正式测量时把显示的 \(A/B\) 除以
\(g\)，即可恢复真实电压幅值比。截距与同轴最大投影的比值不受各自通道增益影响，所以相位读数仍可直接使用。

\begin{enumerate}
\def\labelenumi{\arabic{enumi}.}
\setcounter{enumi}{1}
\item
  若示波器输入电阻为有限值
  \(R_{\rm in}\)，第一次接法中电感支路不再严格开路。画出含
  \(R_{\rm in}\) 的等效电路，说明测得的 \(C\) 会出现系统偏差，并判断
  \(R_{\rm in}\) 越大时结果如何趋近理想值。
\item
  若图 2 只给相位差而不给幅值比，说明为何只能确定
\end{enumerate}

\[
\frac{\omega L}{R+r}=\tan\alpha,
\]

而不能分别求出 \(r,L\)。这一变式强调一个复参数需要模和辐角两条独立信息。

\Needspace{23\baselineskip}

\begin{enumerate}
\def\labelenumi{\arabic{enumi}.}
\setcounter{enumi}{3}
\item
  把工作频率改为
  \(2.00\,\mathrm{kHz}\)，保持元件不变，预测两幅李萨如图的幅值比和相位差如何变化。第一接法更接近高通通带，第二接法的感抗加倍，输出幅值减小、滞后角增大。
\item
  用求得的 \(L,C\) 计算谐振时的感抗与容抗，验证
\end{enumerate}

\[
\omega_0L=\frac1{\omega_0C}.
\]

再说明内阻 \(r\)
会影响串联谐振的峰值和品质因数，却不改变理想条件下的谐振频率公式。

\newpage

\section{第 7 题}\label{ux7b2c-7-ux9898}

海平面空气温度、压强分别为 \(T_0,P_0\)，空气平均摩尔质量为
\(\mu\)，重力加速度为 \(g\)，气体常数为 \(R\)。

（1）若温度不随高度变化，求 \(P(z)\)。

（2）实际压强满足

\[
P(z)=P_0(1-bz)^\lambda,
\]

求温度 \(T(z)\)。

（3）密封热气球体积为 \(V_0\)，海平面起飞时囊内空气温度为
\(2T_0\)。设备等质量（不含囊内空气和配重）为
\(m_0\)；飞行中内外压强相等，能稳定飞行的最大高度为
\(h\)。求起飞时配重质量 \(m\)。

\begin{figure}[H]
\centering
\includegraphics[width=0.72\linewidth,height=\textheight,keepaspectratio,alt={图7-1　大气薄层、竖直坐标与密封热气球的几何示意}]{../figures/question-07.pdf}
\caption*{图7-1　大气薄层、竖直坐标与密封热气球的几何示意}
\end{figure}

\subsection{一、题型定位}\label{ux4e00ux9898ux578bux5b9aux4f4d-6}

大气的压强、温度和密度由静力平衡联系。密封且定体积的气球内气体质量固定，最大高度还须明确配重如何释放以及给定数据是否相容。

\subsection{二、知识点}\label{ux4e8cux77e5ux8bc6ux70b9-6}

\begin{enumerate}
\def\labelenumi{\arabic{enumi}.}
\tightlist
\item
  大气静力平衡 \(dP/dz=-\rho g\)。
\item
  理想气体质量密度 \(\rho=\mu P/(RT)\)。
\item
  已知 \(P(z)\) 可反求 \(T(z)\)。
\item
  浮力等于排开外界空气的重力。
\item
  最大高度通常对应配重全部释放后的临界平衡。
\end{enumerate}

\subsection{三、物理图像}\label{ux4e09ux7269ux7406ux56feux50cf-6}

压强随高度降低，因为上方空气柱重量变小。热气球上升时外界密度下降；为继续保持平衡，需要逐步释放配重。达到最大高度时可释放的配重已经用尽。

\subsection{四、解题入口}\label{ux56dbux89e3ux9898ux5165ux53e3-6}

对厚度 \(dz\)
的水平空气薄层列力平衡，再用理想气体状态方程消去密度。第三问分别对海平面``带全部配重''和高度
\(h\)``配重耗尽''写浮力平衡，两式相减。

\subsection{五、解析}\label{ux4e94ux89e3ux6790-6}

\begin{figure}[H]
\centering
\includegraphics[width=0.88\linewidth,height=\textheight,keepaspectratio,alt={图7-2　大气微层与热气球受力以及状态量随高度的变化}]{../figures/solution-07.pdf}
\caption*{图7-2　大气微层与热气球受力以及状态量随高度的变化}
\end{figure}

图7-2左侧下表面压强向上、上表面压强向下，两力之差平衡薄层重量；右侧分开外界浮力和恒定囊内气体重量，配重的释放条件随后说明。

图7-2左侧分别画出了空气微层和热气球的受力。微层上、下表面的压力差承担该层空气的重力，因此可直接写出式（1）；热气球在最大高度处由浮力与总重力平衡。图中右侧曲线则用于核对第二问：\(P\)
按给定幂律下降，\(T\) 线性下降，故由 \(\rho\propto P/T\)
得到的密度幂次应比压强幂次少 \(1\)。

\needspace{3\baselineskip}
\textbf{（一）等温大气}

空气薄层静力平衡：

\[
dP=-\rho g\,dz. \tag{1}
\]

理想气体给出

\[
\rho=\frac{\mu P}{RT_0}. \tag{2}
\]

代入：

\[
\frac{dP}{P}
=-\frac{\mu g}{RT_0}dz. \tag{3}
\]

从海平面 \(0\) 积分到 \(z\)：

\[
\ln\frac{P(z)}{P_0}
=-\frac{\mu gz}{RT_0}. \tag{4}
\]

故

\[
\boxed{
P(z)=P_0
\exp\left(-\frac{\mu gz}{RT_0}\right)
}. \tag{5}
\]

\needspace{3\baselineskip}
\textbf{（二）由压强剖面反求温度}

一般情况下仍有

\[
\frac{dP}{dz}
=-\frac{\mu P}{RT}g. \tag{6}
\]

解出温度：

\[
T(z)
=-\frac{\mu gP}{R\,dP/dz}. \tag{7}
\]

题给

\[
P=P_0(1-bz)^\lambda. \tag{8}
\]

求导：

\[
\frac{dP}{dz}
=-\lambda bP_0(1-bz)^{\lambda-1}. \tag{9}
\]

代入式（7）：

\[
\boxed{
T(z)=\frac{\mu g}{R\lambda b}(1-bz)
}. \tag{10}
\]

由 \(T(0)=T_0\)，

\[
\boxed{
b=\frac{\mu g}{R\lambda T_0}
}. \tag{11}
\]

所以更直观地

\[
\boxed{
T(z)=T_0(1-bz)
}. \tag{12}
\]

\needspace{3\baselineskip}
\textbf{（三）外界空气密度}

\[
\rho_{\rm out}(z)
=\frac{\mu P(z)}{RT(z)}. \tag{13}
\]

代入式（8）、式（12）：

\[
\boxed{
\rho_{\rm out}(z)
=\rho_0(1-bz)^{\lambda-1},
\qquad
\rho_0=\frac{\mu P_0}{RT_0}
}. \tag{14}
\]

\needspace{3\baselineskip}
\textbf{（四）海平面与最大高度的平衡}

海平面刚起飞时，浮力等于总重力：

\[
\rho_0V_0g
=(m_0+M_{\rm in}+m)g. \tag{15}
\]

其中囊内空气初始质量

\[
M_{\rm in}
=\frac{\mu P_0V_0}{R(2T_0)}
=\frac12\rho_0V_0. \tag{16}
\]

达到最大高度 \(h\)
时，配重已全部释放。气囊密封且体积固定，所以囊内空气质量 \(M_{\rm in}\)
不随高度改变；临界平衡为

\[
\rho_{\rm out}(h)V_0g
=(m_0+M_{\rm in})g. \tag{17}
\]

若直接用式（15）减式（17），可写成

\[
m
=V_0[\rho_0-\rho_{\rm out}(h)]. \tag{18}
\]

但题目给出了设备质量 \(m_0\)，参考答案采用两次平衡式相除，把同一个体积
\(V_0\) 消去：

\[
\frac{m_0+M_{\rm in}+m}{m_0+M_{\rm in}}
=\frac{\rho_0}{\rho_{\rm out}(h)}. \tag{19}
\]

由式（14）

\[
\frac{\rho_0}{\rho_{\rm out}(h)}
=(1-bh)^{1-\lambda}. \tag{20}
\]

故

\[
m=(m_0+M_{\rm in})
\left[(1-bh)^{1-\lambda}-1\right]. \tag{21}
\]

代入囊内空气质量式（16），得到

\[
\boxed{
m=
\left(
m_0+\frac{\mu P_0V_0}{2RT_0}
\right)
\left[
(1-bh)^{1-\lambda}-1
\right]
}. \tag{22}
\]

式（18）和式（22）并不矛盾：利用最大高度平衡式

\[
m_0+M_{\rm in}=\rho_{\rm out}(h)V_0
\]

即可把式（22）化成式（18）。式（22）显式保留了题给的
\(m_0\)，也是与手写参考答案一致的形式。

\needspace{3\baselineskip}
\textbf{补充推导与适用条件}

\[
\frac{dP}{dz}=-\rho g,\qquad
\rho=\frac{\mu P}{RT}.
\]

第一问给定 \(T\)，因此积分求 \(P\)；第二问反过来给定
\(P(z)\)，因此用导数反求
\(T(z)\)。理解这一``已知谁就消去谁''的结构，比记忆指数大气或线性温度公式更重要。

压强幂律

\[
P=P_0(1-bz)^\lambda
\]

与线性温度

\[
T=T_0(1-bz)
\]

并不是两条独立假设。静力平衡强制它们满足

\[
b=\frac{\mu g}{R\lambda T_0}.
\]

进一步得到密度指数 \(\lambda-1\)，因为
\(\rho\propto P/T\)。因此压强、温度和密度的幂次依次为
\(\lambda,1,\lambda-1\)。

第三问必须区分三个量：

\begin{enumerate}
\def\labelenumi{\arabic{enumi}.}
\tightlist
\item
  囊内气体质量 \(M_{\rm in}\)：气囊密封，所以保持不变；
\item
  气囊体积 \(V_0\)：题设明确不变；
\item
  囊内气体温度：只用来确定海平面的初始气体质量，之后不需要假设其随高度如何变化，因为质量已经由密封条件固定。
\end{enumerate}

海平面时内气体温度为 \(2T_0\)，同压同体积下密度是外界的一半，因此

\[
M_{\rm in}=\frac12\rho_0V_0.
\]

这一结果应在列浮力平衡前单独写出，避免把囊内气体也当成设备质量，或在上升后随外界密度重新计算。

配重公式有两种等价物理图像。两式相减：

\[
m=V_0[\rho_0-\rho(h)]
\]

说明抛掉的质量等于浮力承载能力的减少量；两式相除：

\[
\frac{m_0+M_{\rm in}+m}{m_0+M_{\rm in}}
=\frac{\rho_0}{\rho(h)}
\]

说明承载总质量按外界密度比缩放。题目显式给出
\(m_0\)，所以主答案保留相除形式。

极限检查很有用。\(h\to0\) 时 \((1-bh)^{1-\lambda}\to1\)，所以
\(m\to0\)；若通常 \(\lambda>1\) 且 \(0<bh<1\)，则指数
\(1-\lambda<0\)，括号为正，配重质量为正。若算出负值，通常是把密度比倒置了。

\needspace{3\baselineskip}
\textbf{初始平衡、最大高度与稳定性}

海平面初始平衡已经直接给出

\[
m=\frac12\rho_0V_0-m_0.
\]

因此 \(h\)
不是可以独立任选的第五个数据。沿用参考解的``只能释放配重，最高处配重耗尽''模型，还必须同时满足

\[
(1-bh)^{\lambda-1}=\frac12+\frac{m_0}{\rho_0V_0},
\qquad
h=\frac{1-[1/2+m_0/(\rho_0V_0)]^{1/(\lambda-1)}}b.
\]

正配重与正最高高度需要 \(0\le m_0<\rho_0V_0/2\)、\(\lambda>1\)。若题给
\(h\)
不满足此式，参考形式的配重表达式虽然能算出数字，却不能同时满足海平面初态和最高处平衡。

原卷没有明说释放配重，解析把它作为参考解的操作假设写出。若完全不丢配重，气囊质量和体积均固定，在密度单调下降的大气中原先的平衡高度就是海平面，不会另有任意正的最大平衡高度。\(\lambda =1\)
时外界密度恒定，竖直平衡是中性；\(0<\lambda<1\)
时密度反而上升，不能用同一个``配重耗尽的稳定上限''解释。

对固定剩余质量作微小竖直偏移，净力导数为

\[
\frac{dF_z}{dz}=-gV_0\rho_0(\lambda-1)b(1-bz)^{\lambda-2},
\]

在 \(\lambda>1\)
分支为负，才具有回复性。由密封、定体积、内外同压，还可反求囊内温度
\(T_{\rm in}(z)=2T_0(1-bz)^\lambda\)；它不等于始终
\(2T(z)\)。两种温度控制条件属于不同模型。

\needspace{3\baselineskip}
\textbf{不丢配重与丢配重的两条高度关系}

令海平面空气密度为 \(\rho_0\)。海平面囊内温度为
\(2T_0\)、压强与外界相同，故气体质量为
\(\rho_0V_0/2\)。在不改变气体质量和体积时，该质量不会随高度再次变成
\(\rho(z)V_0/2\)。

在仍携带配重 \(m(z)\) 的平衡状态，应有

\[
m(z)=\rho_0V_0(1-bz)^{\lambda-1}-m_0-\frac12\rho_0V_0.
\]

对 \(\lambda>1\) 求导，

\[
\frac{dm}{dz}=-\rho_0V_0(\lambda-1)b(1-bz)^{\lambda-2}<0.
\]

因此若通过逐步释放配重维持平衡，配重必须随高度减少；配重耗尽的高度才是该操作模型下的上限。若完全不释放，外界密度随高度下降而总重不变，向上偏离后浮力小于重力，气球返回海平面。

所有幂律式还要求
\(1-bz>0\)。把最高高度写到温度已经为零或负值的区域，超出了题设大气模型；应在求根前检查
\(0\le z<1/b\)，再检验配重非负。

\subsection{六、关键洞察}\label{ux516dux5173ux952eux6d1eux5bdf-6}

解析图左侧下表面压强向上、上表面压强向下，两力之差平衡薄层重量；右侧分开外界浮力和恒定囊内气体重量，配重的释放条件随后说明。

\subsection{七、易错点}\label{ux4e03ux6613ux9519ux70b9-6}

\begin{enumerate}
\def\labelenumi{\arabic{enumi}.}
\tightlist
\item
  把摩尔质量 \(\mu\) 当成单个分子质量。使用气体常数 \(R\)
  时，物质的量为质量除以摩尔质量，故 \(\rho=\mu P/(RT)\)。
\item
  薄层受力方向写反。下表面压力向上、上表面压力向下，得到
  \(dP/dz=-\rho g\)，压强必须随高度降低。
\item
  等温积分后漏掉边界条件 \(P(0)=P_0\)，或把指数写成正号。
\item
  第二问看到幂律压强仍沿用等温结论。此时应先求
  \(dP/dz\)，再由静力方程解出 \(T\)。
\item
  从
\end{enumerate}

\[
T=-\frac{\mu gP}{R\,dP/dz}
\]

误判温度为负。因为 \(dP/dz<0\)，整体结果为正。 6. 密度幂次写成
\(\lambda+1\)。应由 \(\rho\propto P/T\) 得到 \(\lambda-1\)。 7.
计算囊内空气质量时使用温度 \(T_0\)，漏掉海平面囊内为
\(2T_0\)，从而少了因子 \(1/2\)。 8.
浮力使用囊内空气密度。浮力等于排开的外界空气重量，应使用
\(\rho_{\rm out}(z)V_0g\)。 9.
认为``密封''意味着囊内压强不变。题设说内外压强总平衡，所以囊内压强随高度变化；真正保持不变的是气体物质的量和题设给定的体积。
10. 最大高度平衡式仍保留配重 \(m\)。最大高度对应可释放配重已经耗尽。 11.
两次平衡相除时把密度比写成 \(\rho(h)/\rho_0\)，导致 \(m<0\)。

\subsection{八、课堂讲解建议}\label{ux516bux8bfeux5802ux8bb2ux89e3ux5efaux8bae-6}

讲解可围绕两张守恒表展开：空气薄层的力平衡表和热气球的质量收支表。

\begin{enumerate}
\def\labelenumi{\arabic{enumi}.}
\tightlist
\item
  对面积 \(A\)、厚度 \(dz\) 的空气薄层逐项写力：下表面 \(P(z)A\)
  向上，上表面 \([P(z)+dP]A\) 向下，重力 \(\rho gA\,dz\) 向下。约去
  \(A\) 后再令 \(dz\to0\)。
\item
  第一问把 \(\rho=\mu P/(RT_0)\)
  代入，强调变量可分离；积分上下限同时写出，避免漏常数。
\item
  第二问不积分。把
\end{enumerate}

\[
T=-\frac{\mu gP}{R P'}
\]

看成``由气压计曲线反演温度''的公式，再代入题给幂律。这样学生能理解题目不是重复第一问。
4.
在图7-2旁列三列：外界排开空气质量、设备与气囊内空气、配重。海平面一行含
\(m\)，高度 \(h\) 一行不含 \(m\)；\(m_0,M_{\rm in},V_0\) 两行相同。 5.
先由理想气体算
\(M_{\rm in}\)，再列两次浮力平衡。相减和相除两种路线都展示，但把相除所得、含
\(m_0\) 的形式作为最终答案。 6. 用 \(h=0\)
做即时检验。若学生表达式不能给出
\(m=0\)，无需继续化简，应先检查密度比方向或指数符号。

还应明确模型适用域
\(1-bz>0\)。幂律大气只在该范围内有正温度和正压强；不能把公式外推到
\(z\ge1/b\)。

\subsection{九、变式训练}\label{ux4e5dux53d8ux5f0fux8badux7ec3-6}

\begin{enumerate}
\def\labelenumi{\arabic{enumi}.}
\tightlist
\item
  若气囊柔软、囊内温度恒为 \(2T(z)\)，且气体物质的量 \(n\)
  固定，内外压强相等，则
\end{enumerate}

\[
V(z)=\frac{nR[2T(z)]}{P(z)}.
\]

代入 \(P(z),T(z)\)，求 \(V(z)\) 的高度依赖，再写浮力

\[
F_b(z)=\rho_{\rm out}(z)V(z)g.
\]

比较它与定体积模型的差别。

\begin{enumerate}
\def\labelenumi{\arabic{enumi}.}
\setcounter{enumi}{1}
\tightlist
\item
  若温度随高度线性下降
\end{enumerate}

\[
T=T_0-\alpha z
\]

但不预先给压强，重新积分

\[
\frac{dP}{P}
=-\frac{\mu g\,dz}{R(T_0-\alpha z)}
\]

并把结果写成幂律。由此辨认本题中的 \(\lambda\) 与 \(\alpha\) 的关系。

\begin{enumerate}
\def\labelenumi{\arabic{enumi}.}
\setcounter{enumi}{2}
\tightlist
\item
  若到达高度 \(h\) 时仍剩配重 \(m_r\)，把最大高度平衡式改为
\end{enumerate}

\[
\rho(h)V_0=m_0+M_{\rm in}+m_r,
\]

求已经释放的配重
\(m-m_r\)。说明``给定最大高度''与``给定某一中途高度''在边界条件上的区别。

\begin{enumerate}
\def\labelenumi{\arabic{enumi}.}
\setcounter{enumi}{3}
\tightlist
\item
  取 \(bh\ll1\)，用
\end{enumerate}

\[
(1-bh)^{1-\lambda}
\simeq1+(\lambda-1)bh
\]

求配重质量的一阶近似，并解释其与高度近似成正比。再估计二阶修正何时不可忽略。

\newpage

\section{第 8 题}\label{ux7b2c-8-ux9898}

一束部分椭圆偏振光沿 \(z\) 方向传播。用理想线偏振检偏器测量：透振轴沿
\(x\) 时透射强度最大，为 \(1.5I_0\)；沿 \(y\) 时最小，为 \(I_0\)。

（1）透振轴与 \(x\) 轴成 \(\theta\)
时，求透射强度，并说明是否依赖入射光中非偏振成分的多少。

（2）先让光通过一个四分之一波片，波片快慢轴沿
\(x,y\)，再通过检偏器。发现检偏器透振轴与 \(x\) 成 \(30^\circ\)
时强度最大。求最大强度及入射光中非偏振成分所占比例。

\begin{figure}[H]
\centering
\includegraphics[width=0.76\linewidth,height=\textheight,keepaspectratio,alt={图8-1　部分椭圆偏振光、四分之一波片与旋转检偏器}]{../figures/question-08.pdf}
\caption*{图8-1　部分椭圆偏振光、四分之一波片与旋转检偏器}
\end{figure}

\subsection{一、题型定位}\label{ux4e00ux9898ux578bux5b9aux4f4d-7}

检偏器沿两主轴的读数决定旋转后的强度，却不能单独区分圆偏振和非偏振成分。四分之一波片把相位信息转成方向信息，才能求出非偏振比例。

\subsection{二、知识点}\label{ux4e8cux77e5ux8bc6ux70b9-7}

\begin{enumerate}
\def\labelenumi{\arabic{enumi}.}
\tightlist
\item
  部分偏振光可分为互不相干的非偏振部分和完全偏振部分。
\item
  非偏振光通过任意理想检偏器都透过一半强度。
\item
  椭圆偏振的两个主轴分量相互正交。
\item
  四分之一波片改变相位差，不改变两个轴向分量的强度。
\item
  相位补偿后若偏振部分变为线偏振，透振轴方向由振幅比决定。
\end{enumerate}

\subsection{三、物理图像}\label{ux4e09ux7269ux7406ux56feux50cf-7}

沿 \(x,y\) 测得的两个端点强度都含相同的非偏振贡献
\(I_u/2\)。四分之一波片使椭圆偏振部分的两正交分量同相，从而合成为与
\(x\) 轴成 \(30^\circ\) 的线偏振；非偏振部分仍均匀贡献。

\subsection{四、解题入口}\label{ux56dbux89e3ux9898ux5165ux53e3-7}

设非偏振强度为 \(I_u\)，完全偏振部分沿 \(x,y\) 主轴的强度为
\(I_{ex},I_{ey}\)。先写 \(x,y\) 两个测量方程，再写任意角的投影。第二问由

\[
\tan30^\circ=\sqrt{I_{ey}/I_{ex}}
\]

增加第三个方程。

\subsection{五、解析}\label{ux4e94ux89e3ux6790-7}

\begin{figure}[H]
\centering
\includegraphics[width=0.88\linewidth,height=\textheight,keepaspectratio,alt={图8-2　偏振强度分解、波片相位补偿和最大透射方向}]{../figures/solution-08.pdf}
\caption*{图8-2　偏振强度分解、波片相位补偿和最大透射方向}
\end{figure}

图8-2的蓝、红曲线分别按波片前后强度式绘制；红线最大方向为30度，谷值为
\(0.75I_0\)。谷值对应非偏振分量的一半。

\needspace{3\baselineskip}
\textbf{（一）沿主轴的测量}

非偏振部分通过检偏器的强度恒为 \(I_u/2\)。沿 \(x\)：

\[
\frac{I_u}{2}+I_{ex}=1.5I_0. \tag{1}
\]

沿 \(y\)：

\[
\frac{I_u}{2}+I_{ey}=I_0. \tag{2}
\]

\needspace{3\baselineskip}
\textbf{（二）任意透振轴方向}

透振轴单位矢量为

\[
\hat{\boldsymbol e}_\theta
=\cos\theta\,\hat{\boldsymbol x}
+\sin\theta\,\hat{\boldsymbol y}. \tag{3}
\]

两个相互正交主轴分量分别按马吕斯定律投影：

\[
I(\theta)
=\frac{I_u}{2}
+I_{ex}\cos^2\theta
+I_{ey}\sin^2\theta. \tag{4}
\]

利用式（1）、式（2）：

\[
\begin{aligned}
I(\theta)
&=\left(\frac{I_u}{2}+I_{ex}\right)\cos^2\theta\\
&\quad+
\left(\frac{I_u}{2}+I_{ey}\right)\sin^2\theta\\
&=1.5I_0\cos^2\theta+I_0\sin^2\theta.
\end{aligned} \tag{5}
\]

所以

\[
\boxed{
I(\theta)=I_0\left(1+\frac12\cos^2\theta\right)
}. \tag{6}
\]

该结果只由两个主轴端点强度决定，与 \(I_u\) 的具体数值无关。

\needspace{3\baselineskip}
\textbf{（三）波片后的最大方向}

四分之一波片不改变
\(I_{ex},I_{ey}\)，但把相位差调整为使完全偏振部分成为线偏振。其合成方向满足

\[
\tan30^\circ
=\frac{E_y}{E_x}
=\sqrt{\frac{I_{ey}}{I_{ex}}}. \tag{7}
\]

因此

\[
\frac{I_{ey}}{I_{ex}}=\frac13. \tag{8}
\]

令

\[
I_{ey}=\frac13I_{ex}. \tag{9}
\]

式（1）减式（2）：

\[
I_{ex}-I_{ey}=0.5I_0. \tag{10}
\]

代入式（9）：

\[
\frac23I_{ex}=0.5I_0. \tag{11}
\]

故

\[
\boxed{
I_{ex}=0.75I_0,\qquad
I_{ey}=0.25I_0
}. \tag{12}
\]

再由式（2）：

\[
\frac{I_u}{2}+0.25I_0=I_0, \tag{13}
\]

所以

\[
\boxed{I_u=1.50I_0}. \tag{14}
\]

\needspace{3\baselineskip}
\textbf{（四）最大透射强度与非偏振比例}

完全偏振部分总强度

\[
I_e=I_{ex}+I_{ey}=I_0. \tag{15}
\]

检偏器与其线偏振方向一致时，偏振部分全部透过；非偏振部分仍透过一半。因此

\[
\boxed{
I_{\max}=I_e+\frac{I_u}{2}
=I_0+0.75I_0
=1.75I_0
}. \tag{16}
\]

入射总强度

\[
I_{\rm total}=I_e+I_u=2.50I_0. \tag{17}
\]

非偏振成分比例

\[
\boxed{
\eta_u=\frac{I_u}{I_{\rm total}}
=\frac{1.50}{2.50}
=0.600
}. \tag{18}
\]

即 \(60.0\%\)。

\needspace{3\baselineskip}
\textbf{补充推导与适用条件}

\[
I_{\rm in}=I_u+I_e.
\]

非偏振部分没有固定电场方向，通过任何理想检偏器都贡献
\(I_u/2\)；完全椭圆偏振部分沿主轴 \(x,y\)
的两个正交分量具有确定相位关系，但在只测平均强度时，沿任意线偏振方向的投影可写成

\[
I_{ex}\cos^2\theta+I_{ey}\sin^2\theta.
\]

第一问中出现一个很有用的消元：

\[
\frac{I_u}{2}
=\frac{I_u}{2}(\cos^2\theta+\sin^2\theta).
\]

把它分别并入 \(x,y\)
两项后，任意角强度只由沿主轴测得的两个端点值决定。因此 \(I(\theta)\)
与非偏振成分怎样分配无关。这里的``不依赖''不是说非偏振光没有贡献，而是说它的贡献已经包含在
\(1.5I_0\) 与 \(I_0\) 两个测量值中。

只靠旋转检偏器，方程

\[
\frac{I_u}{2}+I_{ex}=1.5I_0,\qquad
\frac{I_u}{2}+I_{ey}=I_0
\]

有三个未知量，信息不足。四分之一波片提供的不是第三个强度端点，而是相位补偿后的方向信息。椭圆偏振主轴分量原本相差
\(\pm\pi/2\)；波片再引入 \(\mp\pi/2\)
后，两分量同相或反相，合成为线偏振。最大方向给出电场振幅比：

\[
\tan30^\circ=\frac{E_y}{E_x},
\]

而强度正比于振幅平方，所以

\[
\frac{I_{ey}}{I_{ex}}=\tan^230^\circ=\frac13.
\]

必须经过平方，不能把强度比直接写成 \(\tan30^\circ\)。

四分之一波片只改变相位，不改变沿自身快慢轴的振幅，故 \(I_{ex},I_{ey}\)
在波片前后相同。理想波片也不吸收能量，因此入射总强度

\[
I_{\rm total}=I_u+I_{ex}+I_{ey}=2.50I_0
\]

与波片后相同。检偏器最大时完全偏振部分全透过、非偏振部分只透过一半，所以
\(I_{\max}=1.75I_0\) 必小于总强度 \(2.50I_0\)，这是必要的能量检查。

方向 \(30^\circ\)
只确定振幅比的绝对值。若两分量补偿后反相，线偏振方向可写成 \(-30^\circ\)
或
\(150^\circ\)；题图给出的最大方向与波片快慢轴选择共同固定符号。本题求强度与非偏振比例时只需要平方比，方向符号不影响数值结论。

\needspace{3\baselineskip}
\textbf{相干矩阵对交叉项的检查}

归一化到强度单位，初始光束的二阶相干矩阵可写为

\[
J=\begin{pmatrix}1.5I_0&iK\\-iK&I_0\end{pmatrix}.
\]

因为检偏强度在 \(x,y\)
为极值，\(\operatorname{Re}J_{xy}=0\)，所以线性检偏强度中的交叉项确实为零；这来自主轴相位关系，不是``正交分量永远不干涉''。四分之一波片把虚的非对角项变为实项，记符号取正，则

\[
I_{\rm after}(\theta)=1.25I_0+0.25I_0\cos2\theta+K\sin2\theta.
\]

最大方向为 \(30^\circ\)，满足 \(\tan60^\circ=K/(0.25I_0)\)，故
\(K=\sqrt3 I_0/4\)。两个本征强度为 \(1.75I_0\) 和
\(0.75I_0\)。非偏振部分是两轴相同的最小本征值，强度为
\(2\times0.75I_0=1.50I_0\)；余下秩一部分强度
\(I_0\)。由此独立恢复全部结果。

取 \(I_0=1\)，波片前后曲线分别为 \(1.25+0.25\cos2\theta\) 与
\(1.25+0.5\cos[2(\theta-30^\circ)]\)。图8-2的峰值、谷值、周期和最大方向按这两式画出。理想无损波片保持矩阵迹
\(2.5I_0\) 以及本征值，不能增加总光强或偏振度；本题偏振度是
\(I_0/(2.5I_0)=0.4\)。

\subsection{六、关键洞察}\label{ux516dux5173ux952eux6d1eux5bdf-7}

解析图的蓝、红曲线分别按波片前后强度式绘制；红线最大方向为30度，谷值为
\(0.75I_0\)。谷值对应非偏振分量的一半。

\subsection{七、易错点}\label{ux4e03ux6613ux9519ux70b9-7}

\begin{enumerate}
\def\labelenumi{\arabic{enumi}.}
\tightlist
\item
  认为非偏振光通过检偏器的强度随 \(\theta\)
  变化。理想非偏振光在任意方向的平均投影都为总强度的一半。
\item
  把沿 \(x,y\) 测得的 \(1.5I_0,I_0\) 直接当作椭圆偏振分量
  \(I_{ex},I_{ey}\)。两者都还包含同一个 \(I_u/2\)。
\item
  任意角投影中写出交叉干涉项。对椭圆偏振的正交主轴分量，时间平均后的检偏强度可用主轴强度投影相加；非偏振部分又与偏振部分不相干。
\item
  第一问留下未知的 \(I_u\)，没有利用 \(\cos^2\theta+\sin^2\theta=1\)
  消元。
\item
  把强度比直接当作电场振幅比，写成
  \(I_{ey}/I_{ex}=\tan30^\circ\)。正确关系应平方。
\item
  认为四分之一波片改变
  \(I_{ex},I_{ey}\)。理想波片只改变两个轴向分量的相位差，不改变强度。
\item
  默认任意椭圆偏振通过任意四分之一波片都会变成线偏振。只有波片轴与椭圆主轴重合且相位延迟恰好补偿时才成立；题设正是这一特殊配置。
\item
  最大透射强度只取 \(I_e\)，漏加非偏振部分的 \(I_u/2\)。
\item
  入射总强度误写成检偏器最大透射强度。检偏器仍损失一半非偏振光，所以
  \(I_{\rm total}=I_e+I_u\)。
\item
  非偏振比例用 \(I_u/(I_{ex}+I_{ey})\)，分母漏掉非偏振自身。正确分母是
  \(I_u+I_e\)。
\end{enumerate}

\subsection{八、课堂讲解建议}\label{ux516bux8bfeux5802ux8bb2ux89e3ux5efaux8bae-7}

先用检偏强度消去非偏振分量；讨论波片时换用电场振幅，求出方向后再计算透射强度。

\begin{enumerate}
\def\labelenumi{\arabic{enumi}.}
\tightlist
\item
  先画三列表格：非偏振部分、椭圆偏振 \(x\) 分量、椭圆偏振 \(y\) 分量。沿
  \(x\)、沿 \(y\)、任意 \(\theta\) 三行分别填写
  \(I_u/2\)、\(I_{ex}\cos^2\theta\)、\(I_{ey}\sin^2\theta\)。
\item
  在第一问只做强度运算。把 \(I_u/2\) 乘以 \(\cos^2\theta+\sin^2\theta\)
  后分组，现场得到
\end{enumerate}

\[
I(\theta)=I_x\cos^2\theta+I_y\sin^2\theta.
\]

让学生解释为什么无需知道 \(I_u\)。 3.
进入波片部分时切换到电场振幅层。另画一幅振幅合成小图，两箭头长度取
\(\sqrt{I_{ex}}\)、\(\sqrt{I_{ey}}\)，合成方向满足正切关系。图8-2画的是检偏强度曲线，不宜把曲线高度读成电场振幅。
4. 用

\[
\tan30^\circ=\sqrt{I_{ey}/I_{ex}}
\]

求出强度比后，再回到强度层联立式（1）、式（2）。明确每次平方或开方的物理含义。
5. 结尾做能量清单：

\[
I_u=1.50I_0,\quad
I_e=1.00I_0,\quad
I_{\rm total}=2.50I_0.
\]

检偏器最大透过 \(1.00I_0+0.75I_0=1.75I_0\)，最少仍应至少含非偏振的
\(0.75I_0\)。

可让学生把检偏器转到
\(30^\circ+90^\circ=120^\circ\)，计算波片后最小透射强度。完全线偏振部分被挡住，只剩
\(I_u/2=0.75I_0\)，与最大值之差正好等于偏振部分总强度 \(I_e=I_0\)。

\subsection{九、变式训练}\label{ux4e5dux53d8ux5f0fux8badux7ec3-7}

\begin{enumerate}
\def\labelenumi{\arabic{enumi}.}
\tightlist
\item
  波片后把检偏器转到
  \(120^\circ\)，求最小强度。由完全偏振部分与该方向正交，得到
\end{enumerate}

\[
I_{\min,\rm after}=\frac{I_u}{2}=0.75I_0.
\]

验证 \(I_{\max}-I_{\min,\rm after}=I_e\)。

\begin{enumerate}
\def\labelenumi{\arabic{enumi}.}
\setcounter{enumi}{1}
\tightlist
\item
  若检偏器最大方向为一般角 \(\beta\)，则
\end{enumerate}

\[
\frac{I_{ey}}{I_{ex}}=\tan^2\beta.
\]

结合已知的主轴最大、最小透射强度 \(I_x,I_y\)，推导 \(I_u,I_{ex},I_{ey}\)
关于 \(\beta,I_x,I_y\) 的通式，并讨论哪些 \(\beta\)
会给出负强度从而不物理。

\begin{enumerate}
\def\labelenumi{\arabic{enumi}.}
\setcounter{enumi}{2}
\item
  若四分之一波片快慢轴相对 \(x,y\) 旋转
  \(45^\circ\)，说明原椭圆主轴分量不再沿波片本征轴传播，不能直接使用式（7）。可用
  Jones 向量先旋转坐标、施加相位延迟、再旋回。
\item
  若波片相位延迟为未知
  \(\delta\)，仅测检偏器最大、最小强度及其方向，讨论可获得的独立信息数。说明通常还需再改变一次波片方向或增加另一组测量，才能同时反演
  \(I_u\) 与 \(\delta\)。
\end{enumerate}

\Needspace{14\baselineskip}

\begin{enumerate}
\def\labelenumi{\arabic{enumi}.}
\setcounter{enumi}{4}
\tightlist
\item
  用 Stokes 参数表示第一问：
\end{enumerate}

\[
I(\theta)=\frac12\left(S_0+S_1\cos2\theta+S_2\sin2\theta\right).
\]

本题主轴沿 \(x,y\)，所以 \(S_2=0\)。由最大、最小值求
\(S_0,S_1\)，再说明为什么不经过波片无法确定与圆偏振相关的 \(S_3\)。
\end{document}
